% Project 1 — SIMD Advantage Profiling
% Build: `latexmk` in this directory, or `make simd` from the repository root.
\documentclass[11pt]{article}
%% Packages and global settings shared by all ACS reports.
%% Loaded from each project's main.tex:  %% Packages and global settings shared by all ACS reports.
%% Loaded from each project's main.tex:  %% Packages and global settings shared by all ACS reports.
%% Loaded from each project's main.tex:  \input{../../common/latex/preamble}

\usepackage[T1]{fontenc}
\usepackage{lmodern}
\usepackage{microtype}
\usepackage[margin=1in]{geometry}
\usepackage{amsmath,amssymb}
\usepackage{graphicx}
\usepackage{booktabs,tabularx,multirow}
\usepackage{siunitx}
\usepackage{subcaption}
\usepackage[dvipsnames]{xcolor}
\usepackage{tcolorbox}
\tcbuselibrary{breakable}
\usepackage{listings}
\usepackage[shortlabels]{enumitem}
\usepackage{csquotes}

% Plots are drawn from the processed CSVs in <project>/latex/data/ (synced from <project>/code).
\usepackage{pgfplots}
\usepackage{pgfplotstable}
\pgfplotsset{compat=1.18}
\usepgfplotslibrary{groupplots}

\usepackage[backend=bibtex,style=ieee,sorting=nyt]{biblatex} % nyt: avoids a perpetual "undefined references" rerun loop with the bibtex backend
\addbibresource{references.bib}

\usepackage[hidelinks]{hyperref}
\usepackage{xurl}           % allow URLs/paths to break anywhere
\setlength{\emergencystretch}{3em} % let TeX loosen lines instead of overrunning the margin
\usepackage[capitalise,noabbrev]{cleveref}

% ---------------------------------------------------------------- siunitx
\sisetup{
  per-mode = symbol,
  group-digits = integer,
  separate-uncertainty = true,
}
\DeclareSIUnit\flop{FLOP}
\DeclareSIUnit\iops{IOPS}
\DeclareSIUnit\op{op}
\DeclareSIUnit\elem{elem}
\DeclareSIUnit\cycle{cycle}
% Binary prefixes are built in: \qty{128}{\kibi\byte}, \qty{3.2}{\gibi\byte\per\second}

% ---------------------------------------------------------------- listings
\lstset{
  basicstyle = \ttfamily\footnotesize,
  columns = fullflexible,
  keepspaces = true,
  breaklines = true,
  frame = single,
  framerule = 0.3pt,
  rulecolor = \color{black!30},
  showstringspaces = false,
}
% Short disassembly / compiler-report excerpts:  \begin{lstlisting}[style=excerpt] ... \end{lstlisting}
\lstdefinestyle{excerpt}{basicstyle=\ttfamily\scriptsize, frame=leftline}


\usepackage[T1]{fontenc}
\usepackage{lmodern}
\usepackage{microtype}
\usepackage[margin=1in]{geometry}
\usepackage{amsmath,amssymb}
\usepackage{graphicx}
\usepackage{booktabs,tabularx,multirow}
\usepackage{siunitx}
\usepackage{subcaption}
\usepackage[dvipsnames]{xcolor}
\usepackage{tcolorbox}
\tcbuselibrary{breakable}
\usepackage{listings}
\usepackage[shortlabels]{enumitem}
\usepackage{csquotes}

% Plots are drawn from the processed CSVs in <project>/latex/data/ (synced from <project>/code).
\usepackage{pgfplots}
\usepackage{pgfplotstable}
\pgfplotsset{compat=1.18}
\usepgfplotslibrary{groupplots}

\usepackage[backend=bibtex,style=ieee,sorting=nyt]{biblatex} % nyt: avoids a perpetual "undefined references" rerun loop with the bibtex backend
\addbibresource{references.bib}

\usepackage[hidelinks]{hyperref}
\usepackage{xurl}           % allow URLs/paths to break anywhere
\setlength{\emergencystretch}{3em} % let TeX loosen lines instead of overrunning the margin
\usepackage[capitalise,noabbrev]{cleveref}

% ---------------------------------------------------------------- siunitx
\sisetup{
  per-mode = symbol,
  group-digits = integer,
  separate-uncertainty = true,
}
\DeclareSIUnit\flop{FLOP}
\DeclareSIUnit\iops{IOPS}
\DeclareSIUnit\op{op}
\DeclareSIUnit\elem{elem}
\DeclareSIUnit\cycle{cycle}
% Binary prefixes are built in: \qty{128}{\kibi\byte}, \qty{3.2}{\gibi\byte\per\second}

% ---------------------------------------------------------------- listings
\lstset{
  basicstyle = \ttfamily\footnotesize,
  columns = fullflexible,
  keepspaces = true,
  breaklines = true,
  frame = single,
  framerule = 0.3pt,
  rulecolor = \color{black!30},
  showstringspaces = false,
}
% Short disassembly / compiler-report excerpts:  \begin{lstlisting}[style=excerpt] ... \end{lstlisting}
\lstdefinestyle{excerpt}{basicstyle=\ttfamily\scriptsize, frame=leftline}


\usepackage[T1]{fontenc}
\usepackage{lmodern}
\usepackage{microtype}
\usepackage[margin=1in]{geometry}
\usepackage{amsmath,amssymb}
\usepackage{graphicx}
\usepackage{booktabs,tabularx,multirow}
\usepackage{siunitx}
\usepackage{subcaption}
\usepackage[dvipsnames]{xcolor}
\usepackage{tcolorbox}
\tcbuselibrary{breakable}
\usepackage{listings}
\usepackage[shortlabels]{enumitem}
\usepackage{csquotes}

% Plots are drawn from the processed CSVs in <project>/latex/data/ (synced from <project>/code).
\usepackage{pgfplots}
\usepackage{pgfplotstable}
\pgfplotsset{compat=1.18}
\usepgfplotslibrary{groupplots}

\usepackage[backend=bibtex,style=ieee,sorting=nyt]{biblatex} % nyt: avoids a perpetual "undefined references" rerun loop with the bibtex backend
\addbibresource{references.bib}

\usepackage[hidelinks]{hyperref}
\usepackage{xurl}           % allow URLs/paths to break anywhere
\setlength{\emergencystretch}{3em} % let TeX loosen lines instead of overrunning the margin
\usepackage[capitalise,noabbrev]{cleveref}

% ---------------------------------------------------------------- siunitx
\sisetup{
  per-mode = symbol,
  group-digits = integer,
  separate-uncertainty = true,
}
\DeclareSIUnit\flop{FLOP}
\DeclareSIUnit\iops{IOPS}
\DeclareSIUnit\op{op}
\DeclareSIUnit\elem{elem}
\DeclareSIUnit\cycle{cycle}
% Binary prefixes are built in: \qty{128}{\kibi\byte}, \qty{3.2}{\gibi\byte\per\second}

% ---------------------------------------------------------------- listings
\lstset{
  basicstyle = \ttfamily\footnotesize,
  columns = fullflexible,
  keepspaces = true,
  breaklines = true,
  frame = single,
  framerule = 0.3pt,
  rulecolor = \color{black!30},
  showstringspaces = false,
}
% Short disassembly / compiler-report excerpts:  \begin{lstlisting}[style=excerpt] ... \end{lstlisting}
\lstdefinestyle{excerpt}{basicstyle=\ttfamily\scriptsize, frame=leftline}

%% Macros shared by all ACS reports.  %% Macros shared by all ACS reports.  %% Macros shared by all ACS reports.  \input{../../common/latex/macros} after the preamble.

% ================================================================ draft / final
% Draft (default): rubric boxes are shown and missing data is drawn as a placeholder.
% Final (`make final`): rubric boxes are hidden; TODOs and missing data raise LaTeX warnings.
\newif\ifacsdraft
\acsdrafttrue
\ifdefined\acsfinal \acsdraftfalse \fi

% ================================================================ identity
\newcommand{\acsauthor}{Samson Kempiak}
\newcommand{\acsemail}{kempis@rpi.edu}
\newcommand{\acscourse}{ECSE 4320/6320 Advanced Computer Systems}
\newcommand{\acsdue}{October 21, 2026}
\newcommand{\coderepo}{\href{https://github.com/SKempiak/ACS-Code}{\texttt{github.com/SKempiak/ACS-Code}}}

% ================================================================ writing aids
% \todo{...}   something still to write
% \tbd         a number not measured yet
% \code{...}   file names, commands, flags (underscores are fine)
\newcommand{\todo}[1]{%
  \ifacsdraft\else\PackageWarning{acs}{Unresolved TODO}\fi
  \textcolor{red}{\textbf{[TODO:} #1\textbf{]}}}
\newcommand{\tbd}{%
  \ifacsdraft\else\PackageWarning{acs}{Unmeasured value (\string\tbd)}\fi
  \textcolor{orange}{\textbf{??}}}
\newcommand{\code}[1]{\texttt{\detokenize{#1}}}

% Rubric reminder at the top of a section (draft only):  \rubric{Category}{points}{what is graded}
\newtcolorbox{acsrubricbox}{colback=blue!4, colframe=blue!45, boxrule=0.4pt, arc=1pt,
  left=4pt, right=4pt, top=2pt, bottom=2pt, fontupper=\small}
\newcommand{\rubric}[3]{%
  \ifacsdraft\begin{acsrubricbox}\textbf{Rubric --- #1 (#2 pts).} #3\end{acsrubricbox}\fi}

% Pre-experiment hypothesis (common requirement): write and commit it before measuring.
\newtcolorbox{hypothesis}{breakable, colback=black!3, colframe=black!45, boxrule=0.4pt, arc=1pt,
  left=4pt, right=4pt, top=2pt, bottom=2pt, fonttitle=\bfseries\small,
  title={Hypothesis (from first principles, then compared with the data)}}

% Experiment setup table: one "Knob & value \\" row per line.
%   \begin{knobs}
%     Fixed  & kernel = SAXPY, float32 \\
%     Varied & $N = 2^{8} \ldots 2^{26}$ \\
%   \end{knobs}
\newenvironment{knobs}
  {\par\smallskip\noindent\small\tabularx{\linewidth}{@{}>{\bfseries}l X@{}}\toprule}
  {\bottomrule\endtabularx\par\smallskip}

% ================================================================ data & plots
% Processed CSVs live in <project>/latex/data/<name>.csv (common/code/scripts/sync_data.sh).
%
%   \IfDataExists{baseline}{<tikzpicture using data/baseline.csv>}
%
% draws the plot when the file exists, otherwise a placeholder box.
\newcommand{\acsmissing}[1]{%
  \ifacsdraft\else\PackageWarning{acs}{Missing data file #1}\fi
  \fbox{\parbox[c][0.22\textheight][c]{0.9\linewidth}{\centering\small\color{black!55}%
    Data pending: \texttt{\detokenize{#1}}\\[3pt]
    run the experiment in the project's \texttt{code/}, then \texttt{common/code/scripts/sync\_data.sh}}}}
\newcommand{\IfDataExists}[2]{\IfFileExists{data/#1.csv}{#2}{\acsmissing{data/#1.csv}}}

\pgfplotsset{
  % Base style for every report plot.
  acs/.style={
    width=\linewidth, height=0.55\linewidth,
    grid=major, grid style={black!10},
    legend cell align=left, legend style={font=\footnotesize},
    label style={font=\small}, tick label style={font=\footnotesize},
    unbounded coords=discard,
    table/col sep=comma,
  },
  % Keep only rows where column #1 equals #2. Repeat the key to filter on several columns:
  %   \addplot table[x=n, y=ns_per_elem_median, only if={variant}{autovec}, only if={kernel}{saxpy}] {data/baseline.csv};
  /pgfplots/table/only if/.style 2 args={
    /pgfplots/x filter/.append code={%
      \edef\acsFilterValue{\thisrow{#1}}%
      \edef\acsFilterWant{#2}%
      \ifx\acsFilterValue\acsFilterWant\else\def\pgfmathresult{inf}\fi
    },
  },
  % Error bars spanning the interquartile range of a summarized value (tools/summarize columns):
  %   \addplot+[iqr bars] table[x=n, y=ns_per_elem_median, iqr of=ns_per_elem] {...};
  iqr bars/.style={error bars/.cd, y dir=both, y explicit},
  /pgfplots/table/iqr of/.style={
    y error plus expr={\thisrow{#1_q75}-\thisrow{#1_median}},
    y error minus expr={\thisrow{#1_median}-\thisrow{#1_q25}},
  },
}

% Vertical annotation inside an axis, e.g. cache-level boundaries:  \acsvline{32768}{L1d}
\newcommand{\acsvline}[2]{%
  \draw[densely dashed, black!50] ({axis cs:#1,1}|-{rel axis cs:0,0}) -- ({axis cs:#1,1}|-{rel axis cs:0,1})
    node[pos=0.98, anchor=east, rotate=90, font=\scriptsize, fill=white, inner sep=1pt] {#2};}
 after the preamble.

% ================================================================ draft / final
% Draft (default): rubric boxes are shown and missing data is drawn as a placeholder.
% Final (`make final`): rubric boxes are hidden; TODOs and missing data raise LaTeX warnings.
\newif\ifacsdraft
\acsdrafttrue
\ifdefined\acsfinal \acsdraftfalse \fi

% ================================================================ identity
\newcommand{\acsauthor}{Samson Kempiak}
\newcommand{\acsemail}{kempis@rpi.edu}
\newcommand{\acscourse}{ECSE 4320/6320 Advanced Computer Systems}
\newcommand{\acsdue}{October 21, 2026}
\newcommand{\coderepo}{\href{https://github.com/SKempiak/ACS-Code}{\texttt{github.com/SKempiak/ACS-Code}}}

% ================================================================ writing aids
% \todo{...}   something still to write
% \tbd         a number not measured yet
% \code{...}   file names, commands, flags (underscores are fine)
\newcommand{\todo}[1]{%
  \ifacsdraft\else\PackageWarning{acs}{Unresolved TODO}\fi
  \textcolor{red}{\textbf{[TODO:} #1\textbf{]}}}
\newcommand{\tbd}{%
  \ifacsdraft\else\PackageWarning{acs}{Unmeasured value (\string\tbd)}\fi
  \textcolor{orange}{\textbf{??}}}
\newcommand{\code}[1]{\texttt{\detokenize{#1}}}

% Rubric reminder at the top of a section (draft only):  \rubric{Category}{points}{what is graded}
\newtcolorbox{acsrubricbox}{colback=blue!4, colframe=blue!45, boxrule=0.4pt, arc=1pt,
  left=4pt, right=4pt, top=2pt, bottom=2pt, fontupper=\small}
\newcommand{\rubric}[3]{%
  \ifacsdraft\begin{acsrubricbox}\textbf{Rubric --- #1 (#2 pts).} #3\end{acsrubricbox}\fi}

% Pre-experiment hypothesis (common requirement): write and commit it before measuring.
\newtcolorbox{hypothesis}{breakable, colback=black!3, colframe=black!45, boxrule=0.4pt, arc=1pt,
  left=4pt, right=4pt, top=2pt, bottom=2pt, fonttitle=\bfseries\small,
  title={Hypothesis (from first principles, then compared with the data)}}

% Experiment setup table: one "Knob & value \\" row per line.
%   \begin{knobs}
%     Fixed  & kernel = SAXPY, float32 \\
%     Varied & $N = 2^{8} \ldots 2^{26}$ \\
%   \end{knobs}
\newenvironment{knobs}
  {\par\smallskip\noindent\small\tabularx{\linewidth}{@{}>{\bfseries}l X@{}}\toprule}
  {\bottomrule\endtabularx\par\smallskip}

% ================================================================ data & plots
% Processed CSVs live in <project>/latex/data/<name>.csv (common/code/scripts/sync_data.sh).
%
%   \IfDataExists{baseline}{<tikzpicture using data/baseline.csv>}
%
% draws the plot when the file exists, otherwise a placeholder box.
\newcommand{\acsmissing}[1]{%
  \ifacsdraft\else\PackageWarning{acs}{Missing data file #1}\fi
  \fbox{\parbox[c][0.22\textheight][c]{0.9\linewidth}{\centering\small\color{black!55}%
    Data pending: \texttt{\detokenize{#1}}\\[3pt]
    run the experiment in the project's \texttt{code/}, then \texttt{common/code/scripts/sync\_data.sh}}}}
\newcommand{\IfDataExists}[2]{\IfFileExists{data/#1.csv}{#2}{\acsmissing{data/#1.csv}}}

\pgfplotsset{
  % Base style for every report plot.
  acs/.style={
    width=\linewidth, height=0.55\linewidth,
    grid=major, grid style={black!10},
    legend cell align=left, legend style={font=\footnotesize},
    label style={font=\small}, tick label style={font=\footnotesize},
    unbounded coords=discard,
    table/col sep=comma,
  },
  % Keep only rows where column #1 equals #2. Repeat the key to filter on several columns:
  %   \addplot table[x=n, y=ns_per_elem_median, only if={variant}{autovec}, only if={kernel}{saxpy}] {data/baseline.csv};
  /pgfplots/table/only if/.style 2 args={
    /pgfplots/x filter/.append code={%
      \edef\acsFilterValue{\thisrow{#1}}%
      \edef\acsFilterWant{#2}%
      \ifx\acsFilterValue\acsFilterWant\else\def\pgfmathresult{inf}\fi
    },
  },
  % Error bars spanning the interquartile range of a summarized value (tools/summarize columns):
  %   \addplot+[iqr bars] table[x=n, y=ns_per_elem_median, iqr of=ns_per_elem] {...};
  iqr bars/.style={error bars/.cd, y dir=both, y explicit},
  /pgfplots/table/iqr of/.style={
    y error plus expr={\thisrow{#1_q75}-\thisrow{#1_median}},
    y error minus expr={\thisrow{#1_median}-\thisrow{#1_q25}},
  },
}

% Vertical annotation inside an axis, e.g. cache-level boundaries:  \acsvline{32768}{L1d}
\newcommand{\acsvline}[2]{%
  \draw[densely dashed, black!50] ({axis cs:#1,1}|-{rel axis cs:0,0}) -- ({axis cs:#1,1}|-{rel axis cs:0,1})
    node[pos=0.98, anchor=east, rotate=90, font=\scriptsize, fill=white, inner sep=1pt] {#2};}
 after the preamble.

% ================================================================ draft / final
% Draft (default): rubric boxes are shown and missing data is drawn as a placeholder.
% Final (`make final`): rubric boxes are hidden; TODOs and missing data raise LaTeX warnings.
\newif\ifacsdraft
\acsdrafttrue
\ifdefined\acsfinal \acsdraftfalse \fi

% ================================================================ identity
\newcommand{\acsauthor}{Samson Kempiak}
\newcommand{\acsemail}{kempis@rpi.edu}
\newcommand{\acscourse}{ECSE 4320/6320 Advanced Computer Systems}
\newcommand{\acsdue}{October 21, 2026}
\newcommand{\coderepo}{\href{https://github.com/SKempiak/ACS-Code}{\texttt{github.com/SKempiak/ACS-Code}}}

% ================================================================ writing aids
% \todo{...}   something still to write
% \tbd         a number not measured yet
% \code{...}   file names, commands, flags (underscores are fine)
\newcommand{\todo}[1]{%
  \ifacsdraft\else\PackageWarning{acs}{Unresolved TODO}\fi
  \textcolor{red}{\textbf{[TODO:} #1\textbf{]}}}
\newcommand{\tbd}{%
  \ifacsdraft\else\PackageWarning{acs}{Unmeasured value (\string\tbd)}\fi
  \textcolor{orange}{\textbf{??}}}
\newcommand{\code}[1]{\texttt{\detokenize{#1}}}

% Rubric reminder at the top of a section (draft only):  \rubric{Category}{points}{what is graded}
\newtcolorbox{acsrubricbox}{colback=blue!4, colframe=blue!45, boxrule=0.4pt, arc=1pt,
  left=4pt, right=4pt, top=2pt, bottom=2pt, fontupper=\small}
\newcommand{\rubric}[3]{%
  \ifacsdraft\begin{acsrubricbox}\textbf{Rubric --- #1 (#2 pts).} #3\end{acsrubricbox}\fi}

% Pre-experiment hypothesis (common requirement): write and commit it before measuring.
\newtcolorbox{hypothesis}{breakable, colback=black!3, colframe=black!45, boxrule=0.4pt, arc=1pt,
  left=4pt, right=4pt, top=2pt, bottom=2pt, fonttitle=\bfseries\small,
  title={Hypothesis (from first principles, then compared with the data)}}

% Experiment setup table: one "Knob & value \\" row per line.
%   \begin{knobs}
%     Fixed  & kernel = SAXPY, float32 \\
%     Varied & $N = 2^{8} \ldots 2^{26}$ \\
%   \end{knobs}
\newenvironment{knobs}
  {\par\smallskip\noindent\small\tabularx{\linewidth}{@{}>{\bfseries}l X@{}}\toprule}
  {\bottomrule\endtabularx\par\smallskip}

% ================================================================ data & plots
% Processed CSVs live in <project>/latex/data/<name>.csv (common/code/scripts/sync_data.sh).
%
%   \IfDataExists{baseline}{<tikzpicture using data/baseline.csv>}
%
% draws the plot when the file exists, otherwise a placeholder box.
\newcommand{\acsmissing}[1]{%
  \ifacsdraft\else\PackageWarning{acs}{Missing data file #1}\fi
  \fbox{\parbox[c][0.22\textheight][c]{0.9\linewidth}{\centering\small\color{black!55}%
    Data pending: \texttt{\detokenize{#1}}\\[3pt]
    run the experiment in the project's \texttt{code/}, then \texttt{common/code/scripts/sync\_data.sh}}}}
\newcommand{\IfDataExists}[2]{\IfFileExists{data/#1.csv}{#2}{\acsmissing{data/#1.csv}}}

\pgfplotsset{
  % Base style for every report plot.
  acs/.style={
    width=\linewidth, height=0.55\linewidth,
    grid=major, grid style={black!10},
    legend cell align=left, legend style={font=\footnotesize},
    label style={font=\small}, tick label style={font=\footnotesize},
    unbounded coords=discard,
    table/col sep=comma,
  },
  % Keep only rows where column #1 equals #2. Repeat the key to filter on several columns:
  %   \addplot table[x=n, y=ns_per_elem_median, only if={variant}{autovec}, only if={kernel}{saxpy}] {data/baseline.csv};
  /pgfplots/table/only if/.style 2 args={
    /pgfplots/x filter/.append code={%
      \edef\acsFilterValue{\thisrow{#1}}%
      \edef\acsFilterWant{#2}%
      \ifx\acsFilterValue\acsFilterWant\else\def\pgfmathresult{inf}\fi
    },
  },
  % Error bars spanning the interquartile range of a summarized value (tools/summarize columns):
  %   \addplot+[iqr bars] table[x=n, y=ns_per_elem_median, iqr of=ns_per_elem] {...};
  iqr bars/.style={error bars/.cd, y dir=both, y explicit},
  /pgfplots/table/iqr of/.style={
    y error plus expr={\thisrow{#1_q75}-\thisrow{#1_median}},
    y error minus expr={\thisrow{#1_median}-\thisrow{#1_q25}},
  },
}

% Vertical annotation inside an axis, e.g. cache-level boundaries:  \acsvline{32768}{L1d}
\newcommand{\acsvline}[2]{%
  \draw[densely dashed, black!50] ({axis cs:#1,1}|-{rel axis cs:0,0}) -- ({axis cs:#1,1}|-{rel axis cs:0,1})
    node[pos=0.98, anchor=east, rotate=90, font=\scriptsize, fill=white, inner sep=1pt] {#2};}

%% Measurement environment macros, used by common/latex/environment_table.tex.
%% GENERATED by common/code/scripts/collect_env.sh --tex on the test machine.
%% Regenerate with (from the repository root):
%%   common/code/scripts/collect_env.sh --tex > common/latex/system_info.tex
\newcommand{\sysCPU}{Intel(R) Core(TM) i9-10900F CPU @ 2.80GHz}
\newcommand{\sysArch}{x86\_64}
\newcommand{\sysISA}{sse sse2 ssse3 sse4\_1 sse4\_2 avx avx2 fma f16c -- widest 256-bit (ymm)}
\newcommand{\sysCores}{1 socket(s), 10 cores/socket, 1 threads/core (10 of 20 logical CPUs online)}
\newcommand{\sysSMT}{off}
\newcommand{\sysCaches}{L1d 320 KiB (10 instances); L2 2.5 MiB (10 instances); L3 20 MiB (1 instance)}
\newcommand{\sysNUMA}{1 node(s)}
\newcommand{\sysGovernor}{performance}
\newcommand{\sysDriver}{intel\_pstate}
\newcommand{\sysTurbo}{disabled}
\newcommand{\sysEPP}{performance}
\newcommand{\sysIsolation}{no isolcpus}
\newcommand{\sysMemory}{63 GiB total}
\newcommand{\sysDIMM}{2 $\times$ 32 GiB DDR4-3200, one per channel (A, B), running at 2933 MT/s (dmidecode)}
\newcommand{\sysPageSize}{4096 bytes}
\newcommand{\sysHugePages}{always}
\newcommand{\sysOS}{Arch Linux}
\newcommand{\sysKernel}{7.2.9-arch1-1}
\newcommand{\sysCC}{cc (GCC) 16.2.1 20260810}
\newcommand{\sysGCC}{gcc (GCC) 16.2.1 20260810}
\newcommand{\sysClang}{clang version 23.1.3}
\newcommand{\sysPerf}{perf version 7.2.11-1}
\newcommand{\sysPerfParanoid}{0}
\newcommand{\sysFio}{fio-3.42}
\newcommand{\sysMLC}{present}
\newcommand{\sysGPU}{NVIDIA GeForce RTX 3070}
\newcommand{\sysCUDA}{not available (no CUDA)}
\newcommand{\sysCollected}{2026-10-21 20:46 EDT}


% Typewriter text: keep "--" as two hyphens (command-line flags such as --info,
% --cpu, --selftest) instead of the T1 en-dash ligature.
\DisableLigatures[-]{encoding = T1, family = tt*}

% With the bibtex fall-back backend (no biber on the build host), biblatex 3.21 never records the
% previous run's citation list, so the citation-order check of the ieee style's sorting=none
% template fails on every run and reports "There were undefined references" / "Please rerun
% BibTeX". This report cites a single source, so name-year sorting gives the same list and avoids
% the spurious warning.
\ExecuteBibliographyOptions{sorting=nyt}

% Hypothesis boxes: the shared title claims "written before measuring", which git history
% cannot prove for this report; state what was actually done.
\renewtcolorbox{hypothesis}{breakable, colback=black!3, colframe=black!45, boxrule=0.4pt, arc=1pt,
  left=4pt, right=4pt, top=2pt, bottom=2pt, fonttitle=\bfseries\small,
  title={Hypothesis (from first principles, then compared with the data)}}

% Error bars for speedup columns: speedup_q25/_q75 are appended by
% SIMD_Advantage_Profiling/code/scripts/speedup_iqr.awk.
\pgfplotsset{
  % Shorter single-panel plots than the shared default, to keep the report compact.
  acs/.append style={height=0.42\linewidth},
  % Fixed per-variant styles for plots where a variant may still be missing from the data;
  % the legend is drawn with \acsvariantlegend so it never shifts onto the wrong curve.
  var scalar/.style={color=blue, mark=*, mark options={fill=blue}},
  var autovec/.style={color=red, mark=square*, mark options={fill=red}},
  var native/.style={color=brown!80!black, mark=triangle*, mark options={fill=brown!80!black}},
  /pgfplots/table/speedup iqr/.style={
    y error plus expr={\thisrow{speedup_q75}-\thisrow{speedup_vs_scalar}},
    y error minus expr={\thisrow{speedup_vs_scalar}-\thisrow{speedup_q25}},
  },
}
\newcommand{\acsvariantlegend}{%
  \addlegendimage{var scalar}\addlegendentry{scalar}%
  \addlegendimage{var autovec}\addlegendentry{autovec}%
  \addlegendimage{var native}\addlegendentry{native}}

% Let floats share pages with text instead of forming half-empty float pages.
\renewcommand{\topfraction}{0.9}
\renewcommand{\bottomfraction}{0.8}
\renewcommand{\textfraction}{0.07}
\renewcommand{\floatpagefraction}{0.75}
\setcounter{topnumber}{3}
\setlength{\textfloatsep}{10pt plus 2pt minus 4pt}
\setlength{\floatsep}{8pt plus 2pt minus 2pt}
\setlength{\intextsep}{8pt plus 2pt minus 2pt}
\setcounter{totalnumber}{4}

\title{Project 1: SIMD Advantage Profiling}
\author{\acsauthor\\[2pt]\small\acscourse\\\small\href{mailto:\acsemail}{\acsemail}}
\date{\acsdue}

\begin{document}
\maketitle

\begin{abstract}
  We measure scalar-versus-SIMD speedup for four single-threaded kernels (SAXPY, dot product,
  elementwise multiply and a 3-point stencil) on an Intel Core i9-10900F (AVX2+FMA,
  \qty{256}{bit} vectors) at a fixed \qty{2.8}{\giga\hertz}. With cache-resident data a
  \code{-march=native} build reaches about $5.8\times$ over scalar for SAXPY and the generic SSE
  build about $3\times$. In DRAM the gain compresses to $1.88\text{--}2.02\times$ for SAXPY: the
  native kernels sit on the measured bandwidth roof, and what remains is the scalar loop's own
  instruction cost. Float64 usually keeps about half of float32's cache-resident gain. Any non-unit
  stride or gather erases the advantage; the AVX2 hardware gather is even $1.5\text{--}1.6\times$
  slower than scalar in cache. A cache-aware roofline with separately
  measured ceilings (\qty{89.2}{\giga\flop\per\second} FMA peak, triad bandwidth per level) places L3-
  and DRAM-resident native kernels at \SIrange{95}{107}{\percent} of their roof. Our negative cases
  are memory-bound SAXPY and the non-reassociated dot product ($\approx1\times$).
\end{abstract}

\ifacsdraft{\setcounter{tocdepth}{1}\tableofcontents}\fi

% !TeX root = ../main.tex
\section{Introduction}\label{sec:intro}

Single-instruction-multiple-data (SIMD) units let one instruction operate on a whole vector of
elements. On the Intel Core i9-10900F studied here the widest unit is AVX2 with FMA: a
\qty{256}{bit} register holds \num{8} float32 or \num{4} float64 lanes, so the \emph{ideal} speedup
over scalar code is $8\times$ or $4\times$. Loads and stores, loop tails, reduction dependencies
and, above all, memory bandwidth keep real kernels below that. We measure where the benefit
appears, how far it falls short of the lane count, and when it disappears:
\begin{enumerate}
  \item How much faster are auto-vectorized builds than scalar builds for each kernel, and how does
        that change from cache-resident to DRAM-resident working sets? (\cref{sec:baseline,sec:locality})
  \item What do alignment, loop tails, stride and gather access, and data type cost? (\cref{sec:alignment,sec:stride,sec:dtype})
  \item Where do these kernels sit in a roofline model of this machine? (\cref{sec:roofline})
  \item When does the expected SIMD benefit fail to appear? (\cref{sec:negative})
\end{enumerate}

% !TeX root = ../main.tex
\section{Methodology}\label{sec:method}
\rubric{Baseline, correctness, reproducibility}{15}{Reference validation, controlled timing,
reproducible scripts and a documented environment (with \cref{sec:correctness,sec:baseline}).}

\subsection{Platform}
\newcommand{\acspinning}{Each run is pinned to one physical core
(\code{--cpu 2}, \code{sched_setaffinity}); SMT is disabled, so the core has no sibling
thread, and no other pinning/isolation is used.}
%% Measurement environment table (common requirement: measurement discipline).
%% Values come from common/system_info.tex. Before \input-ing this file a project may define
%%   \acspinning   how runs were pinned/isolated
%%   \acsenvextra  extra rows, e.g.  \newcommand{\acsenvextra}{SSD & model, firmware, PCIe link \\}
\providecommand{\acspinning}{\todo{e.g.\ \code{taskset -c 2} on a P-core; isolcpus if used}}
\providecommand{\acsenvextra}{}
\begin{table}[tbp]
  \centering\small
  \caption{Measurement environment.}\label{tab:environment}
  \begin{tabularx}{\linewidth}{@{}>{\raggedright\arraybackslash}p{0.2\linewidth}X@{}}
    \toprule
    CPU               & \sysCPU{} (\sysArch) \\
    SIMD ISA          & \sysISA \\
    Topology          & \sysCores; NUMA: \sysNUMA \\
    SMT               & \sysSMT \\
    Caches            & \sysCaches \\
    Memory            & \sysMemory; DIMMs: \sysDIMM \\
    Pages             & base \sysPageSize; THP: \sysHugePages \\
    OS / kernel       & \sysOS; \sysKernel \\
    Frequency policy  & governor: \sysGovernor; driver: \sysDriver; turbo: \sysTurbo; EPP: \sysEPP \\
    Pinning / isolation & \acspinning; \sysIsolation \\
    Compilers         & \sysGCC; \sysClang \\
    Tools             & perf: \sysPerf{} (perf\_event\_paranoid \sysPerfParanoid); FIO: \sysFio; MLC: \sysMLC \\
    GPU (optional)    & \sysGPU; CUDA \sysCUDA \\
    \acsenvextra
    Env.\ snapshot     & \sysCollected{} (each raw run directory also stores its own \code{env.txt}) \\
    \bottomrule
  \end{tabularx}
\end{table}


This Comet Lake core implements AVX2 and FMA but \emph{not} AVX-512, so the widest vector is
\qty{256}{bit} (\code{ymm}). \code{-march=native} targets AVX2+FMA; the generic \code{x86-64}
target only guarantees SSE2, so the plain auto-vectorized build emits \qty{128}{bit} (\code{xmm})
code. That gap is why \code{autovec} and \code{native} are separate variants (\cref{tab:variants}).

\subsection{Kernels}
\begin{table}[tbp]
  \centering\small
  \caption{Kernels and floating-point operations (FLOPs) per processed element.}\label{tab:kernels}
  \begin{tabular}{@{}lll@{}}
    \toprule
    Kernel & Definition & FLOPs / element \\
    \midrule
    SAXPY               & $y_i \leftarrow a\,x_i + y_i$                         & 2 \\
    Dot product         & $s \leftarrow \sum_i x_i\, y_i$                        & 2 \\
    Elementwise multiply & $z_i \leftarrow x_i\, y_i$                            & 1 \\
    3-point stencil     & $y_i \leftarrow a\,x_{i-1} + b\,x_i + c\,x_{i+1}$      & 5 \\
    \bottomrule
  \end{tabular}
\end{table}
A multiply and an add count as one FLOP each, so an FMA is two. All four kernels are implemented
and validated; SAXPY, with the simplest unit-stride access, is the representative kernel for the
locality, alignment and stride studies. Inputs are a fixed-seed linear-congruential sequence in
$[-0.5, 0.5)$ (no denormals, NaNs or overflow); the gather uses a Fisher--Yates permutation of
$[0,N)$. Float32 SAXPY touches \qty{8}{\byte} per element, so its working set crosses the
\qty{32}{\kibi\byte} L1d near $N\approx\num{4096}$, the \qty{256}{\kibi\byte} L2 near
$N\approx\num{32768}$ and the \qty{20}{\mebi\byte} L3 near $N\approx\num{2.6e6}$; plots mark these
boundaries.

\subsection{Build variants}
\begin{table}[tbp]
  \centering\small
  \caption{Build variants. Only the kernel translation unit (\code{src/kernels.c}) differs; the
  timing harness (\code{main.c}, \code{common.c}) and the reference (\code{reference.c}) are
  compiled identically across variants.}\label{tab:variants}
  \begin{tabularx}{\linewidth}{@{}lX@{}}
    \toprule
    Variant & Kernel compiler flags (from \code{simd_bench --info}) \\
    \midrule
    scalar    & \code{-O3 -fno-tree-vectorize -fno-tree-slp-vectorize} \\
    autovec   & \code{-O3} (generic \code{x86-64} $\Rightarrow$ SSE2, \qty{128}{bit}) \\
    native    & \code{-O3 -march=native} (AVX2+FMA, \qty{256}{bit}) \\
    native512 & not available on this CPU (no AVX-512; \qty{256}{bit} is the widest ISA) \\
    \bottomrule
  \end{tabularx}
\end{table}
The compiler is GCC \sysGCC{}. The reference is built with
\code{-O2 -fno-tree-vectorize -ffp-contract=off}, so it evaluates each expression in strict,
non-fused order. No variant uses \code{-ffast-math}: results stay reproducible, but the compiler
cannot reassociate the dot-product reduction, an effect we measure rather than hide
(\cref{sec:vectorization,sec:negative}).

\subsection{Data collection, timing and hypotheses}
All measurements were taken on \textbf{2026-10-21} on an idle machine, settings checked before each
run and jobs serialized through a lock file: \code{baseline} at 20:38, the other six experiments
from 20:47 to 20:54 (\cref{tab:raw-data}). Beforehand the user masked an ACPI event storm
(GPE~0x42) that had stolen $\approx\SI{13}{\percent}$ of wall time in the companion Cache \& Memory
project; the mask does not survive a reboot (\cref{app:repro}). Governor \sysGovernor{}, turbo
\sysTurbo{} and SMT off (\cref{tab:environment}) fix the core at \qty{2.8}{\giga\hertz}. Two checks
guard against a shared core: no configuration varies by more than $1.36\times$ between its fastest
and slowest repetition (all but one $\le1.32\times$), and the user-space cycle counter, where stolen
time would show as a lower clock, reads \qtyrange{2.78}{2.79}{\giga\hertz} for all \num{96} locality
and \num{123} ceilings configurations.
\begin{itemize}
  \item Timer: \code{CLOCK_MONOTONIC_RAW} (\qty{1}{\nano\second} resolution); each repetition times
        a calibrated batch of at least \qty{25}{\milli\second}, so resolution and call overhead are
        negligible. 3 warm-up invocations, 12 repetitions; median with interquartile-range (IQR)
        error bars. Speedup error bars span
        $[\,q_{25}(\text{SIMD})/q_{75}(\text{scalar}),\; q_{75}(\text{SIMD})/q_{25}(\text{scalar})\,]$
        (\code{scripts/speedup_iqr.awk}).
  \item Results pass through a \code{volatile} sink and kernels live in their own translation unit,
        so no loop is dead-code eliminated. The sweep (\code{scripts/run.sh}) is not shuffled.
  \item An earlier turbo-on sweep (2026-10-12, commit \code{73f8fb8}, raw directory
        \code{20261012-182813}) is not used for any result. Its cache-resident rates were
        $1.72\text{--}1.75\times$ the fixed-clock values (dot product
        \qtyrange{2.43}{2.49}{\giga\flop\per\second} at $0.5$~FLOP/cycle, i.e.\
        $\approx\qty{4.9}{\giga\hertz}$), while DRAM bandwidth does not scale, so its DRAM-resident
        SAXPY speedups were only $1.32\text{--}1.44\times$, the single turbo figure we quote.
\end{itemize}
Each experiment opens with a hypothesis from first principles (lane counts, cache sizes, port
counts). Some were drafted after the superseded turbo sweep, so we treat them as predictions and
state where the data contradicts them.

\subsection{Machine ceilings and hardware counters}\label{sec:ceilings-method}
The roofline (\cref{sec:roofline}) needs ceilings \emph{not} derived from the kernels it explains;
a separate program (\code{src/ceilings.c}, \code{run.sh ceilings}) measures them:
\begin{itemize}
  \item \emph{Peak FLOP rate}: \num{12} independent register accumulators $a \leftarrow a\,m + c$
        (with \num{4}-cycle latency and two FMA ports, $4\times2=8$ chains saturate the ports), as
        scalar \code{mulss}+\code{addss} (\code{scalar}'s ceiling), scalar FMA, SSE
        \code{mulps}+\code{addps} (\code{autovec}, no FMA) and AVX2 \code{vfmadd231ps}/\code{pd}
        (\code{native}).
  \item \emph{Bandwidth}: STREAM-style \emph{read}, \emph{copy} and \emph{triad}
        ($a_i = b_i + s\,c_i$) with AVX2 loads/stores, at a footprint inside each level (L1
        \qty{16}{\kibi\byte}, L2 \qty{128}{\kibi\byte}, L3 \qty{8}{\mebi\byte}, DRAM
        \qty{512}{\mebi\byte}) and over a \qty{4}{\kibi\byte}--\qty{512}{\mebi\byte} sweep. Roofs and
        kernel intensities count actual traffic: explicit loads and stores in L1, where stores hit,
        plus the write-allocate read in L2, L3 and DRAM (\cref{tab:roofline-inputs}). The read
        stream's add chain caps it at one \qty{32}{\byte} load per cycle
        (\qty{89}{\gibi\byte\per\second}), so it is never a roof.
\end{itemize}
\emph{Hardware counters.} The AI agent that ran the benchmarks had no root; the user applied the
root-only settings (governor, turbo, SMT, GPE mask, \code{perf_event_paranoid}~$=$~\sysPerfParanoid{}).
\code{src/perfctr.c} opens the thread's own user-space core-PMU events (cycles, instructions, LLC
misses, L1D read misses) with \code{perf_event_open(2)}, unprivileged up to
\code{perf_event_paranoid}~$=2$, enabled only around each timed repetition (\code{NA} if
unavailable).

% !TeX root = ../main.tex
\section{Correctness Validation}\label{sec:correctness}

Every build variant is checked two independent ways before its timings are trusted
(\code{make test}, \code{scripts/validate.sh}). \emph{Reference comparison:} \code{src/reference.c}
holds plain scalar loops compiled with \code{-fno-tree-vectorize -ffp-contract=off} (strict,
non-fused, left-to-right order); the driver runs kernel and reference on identical inputs and
compares them element by element (or the scalar result, for the dot product) before timing.

\begin{knobs}
  Cases  & kernels $\times$ \{float32, float64\}; $N$ around multiples of the lane count
           ($1,2,3,7,8,9,15,16,17,31,33,64,127,1000,1024,4095,4096,4097,65536$); start offsets
           \numlist{4;8;12;16;20;28;32;60}~\unit{\byte}; strides $2,4,8,16$; gather; explicit AVX2 gather (native only,
           $N \in \{7,8,9,16,17,1000,4096,65536\}$); int32 multiply \\
  Tolerance & relative error $\le 10^{-5}$ (float32) and $\le 10^{-12}$ (float64): loose enough for
             FMA contraction in the vectorized builds, tight enough to catch real bugs \\
  Command & \code{make test} \\
\end{knobs}

\emph{Known-answer self-test} (\code{simd_bench --selftest}): fixed inputs with known results for
every kernel and the AVX2 gather (e.g.\ $y=2x+y$ with $x=1,y=3$ gives $5$), which guards against
kernel and reference sharing a mistake. Remainder sizes and misaligned starts exercise the epilogue
and prologue that later experiments stress.

Result (\code{data/raw/validation/20261021-231629/make_test.txt}): all \num{578} cases pass
(\num{190} each for scalar and autovec, \num{198} for native, which adds the AVX2 gather). The
largest relative error is \num{2.1e-7} (native float32 dot product, $N=127$); every other case,
including all float64 and int32 cases, matches bit for bit. The test coefficients ($a=0.5$; stencil
$0.25, 0.5, 0.25$) make every product exact, so only the reduction exercises the tolerance.

% !TeX root = ../main.tex
\section{Vectorization Verification}\label{sec:vectorization}
\rubric{Vectorization verification}{15}{Clear evidence that SIMD instructions are used, with a correct
reading of vector width, FMA and masking.}

\begin{hypothesis}
  The \code{scalar} build should emit only scalar SSE (\code{ss}/\code{sd}) instructions and report
  zero vectorized loops. \code{autovec}, limited to the generic SSE2 target, should vectorize the
  streaming kernels at \qty{128}{bit} (\num{4} float32 lanes). \code{native} should vectorize them at
  \qty{256}{bit} (\code{ymm}, \num{8} float32 / \num{4} float64 lanes) and fuse the multiply--add in
  SAXPY and the stencil into \code{vfmadd}. We expect the compiler to \emph{refuse} to vectorize the
  dot-product reduction usefully, because summation reassociation is illegal without
  \code{-ffast-math}.
\end{hypothesis}

\Cref{tab:vec-summary} counts, per kernel function and variant, the instructions that use vector
registers or packed operations (\code{objdump}, \code{make report}). The few matches in the scalar
column are register-clearing \code{pxor}/\code{movaps}; GCC reports \emph{zero} vectorized loops
for \code{scalar}.

\begin{table}[tbp]
  \centering\small
  \caption{SIMD instructions per kernel function and build variant.}\label{tab:vec-summary}
  \IfDataExists{vectorization}{%
    \pgfplotstableset{acsvec/.style={
      col sep=comma,
      columns={variant,function,insns,simd_insns,widest_register},
      columns/variant/.style={column name=Build, string type},
      columns/function/.style={column name=Function, string type, string replace*={_}{\_}},
      columns/insns/.style={column name=Insns},
      columns/simd_insns/.style={column name=SIMD},
      columns/widest_register/.style={column name=Widest, string type},
      every head row/.style={before row=\toprule, after row=\midrule},
      every last row/.style={after row=\bottomrule}}}%
    \scriptsize\setlength{\tabcolsep}{2.5pt}%
    \pgfplotstabletypeset[acsvec, skip rows between index={12}{36}]{data/vectorization.csv}\hfill
    \pgfplotstabletypeset[acsvec, skip rows between index={0}{12}, skip rows between index={24}{36}]{data/vectorization.csv}\hfill
    \pgfplotstabletypeset[acsvec, skip rows between index={0}{24}]{data/vectorization.csv}%
  }
\end{table}

\subsection{Compiler report}
\code{make report} (\code{data/raw/vectorization/20261021-001032}) recompiles the unchanged
\code{src/kernels.c} with each variant's flags; \code{objdump} confirms the kernels match the timed
binaries instruction for instruction. \code{-fopt-info-vec} for SAXPY (\code{kernels.c:29}):
\begin{lstlisting}[style=excerpt]
# native  (-O3 -march=native):
src/kernels.c:29:26: optimized: loop vectorized using 32 byte vectors and unroll factor 8
src/kernels.c:29:26: optimized: epilogue loop vectorized using 16 byte vectors and unroll factor 4
# autovec (-O3, generic x86-64):
src/kernels.c:29:26: optimized: loop vectorized using 16 byte vectors and unroll factor 4
src/kernels.c:29:26: optimized: epilogue loop vectorized using 8 byte vectors and unroll factor 2
# scalar  (-O3 -fno-tree-vectorize): 0 loops vectorized
\end{lstlisting}
\qty{32}{\byte} vectors are \num{8} float32 lanes (native), \qty{16}{\byte} are \num{4} (autovec);
GCC's ``unroll factor'' is the vectorization factor, not extra unrolling. The native epilogue is one
\qty{128}{bit} step (if \num{4} or more elements remain) plus up to three scalar iterations;
autovec's is one \qty{8}{\byte} step plus at most one scalar iteration. The loop is
\emph{versioned} on a run-time aliasing check between \code{x} and \code{y}, and GCC does not peel
iterations to reach an aligned start.

\subsection{Disassembly}
The native SAXPY inner loop is one fused multiply--add on \code{ymm} registers:
\begin{lstlisting}[style=excerpt]
vmovups   (%rdi,%rax,1),%ymm0           # load 8 x[i]
vfmadd213ps (%rsi,%rax,1),%ymm1,%ymm0   # ymm0 = a*x + y   (one FMA, 8 lanes)
vmovups   %ymm0,(%rsi,%rax,1)           # store 8 y[i]
\end{lstlisting}
The scalar build handles one element at a time and cannot fuse, because the generic \code{x86-64}
target lacks FMA:
\begin{lstlisting}[style=excerpt]
movss (%rdi,%rax,4),%xmm0  # x[i]
mulss %xmm1,%xmm0          # a*x[i]
addss (%rsi,%rax,4),%xmm0  # + y[i]   (separate mul + add: generic x86-64 has no FMA)
movss %xmm0,(%rsi,%rax,4)  # y[i]
\end{lstlisting}

The dot product is the exception: native multiplies eight lanes at once but sums them one lane at
a time, in source order:
\begin{lstlisting}[style=excerpt]
vmulps (%rdi,%rax,1),%ymm1,%ymm1  # 8 products at once ...
vaddss %xmm1,%xmm0,%xmm0          # ... but summed one lane at a time,
vaddss %xmm3,%xmm0,%xmm0          #     in strict source order
...
\end{lstlisting}

\subsection{Stride and gather code}
\code{saxpy_f32_stride} shows \code{ymm} instructions in \code{native} only because GCC versions
the loop on \code{s == 1}; for $s>1$ it runs a scalar \code{vfmadd213ss} loop (\code{mulss}/\code{addss}
in the other builds). No build vectorizes the plain gather (\enquote{no vectype for stmt}), so we
added an explicit AVX2 version, \code{saxpy_f32_gather_avx2}, compiled only into \code{native}:
\begin{lstlisting}[style=excerpt]
vmovdqu    (%rax),%ymm1                 # 8 indices idx[i..i+7]
vgatherdps %ymm5,(%rdi,%ymm1,4),%ymm0   # gather 8 x[idx]
vgatherdps %ymm7,(%rsi,%ymm1,4),%ymm3   # gather 8 y[idx]
vfmadd132ps %ymm4,%ymm3,%ymm0           # one 8-lane FMA
vmovss     %xmm0,(%rsi,%rdx,4)          # ... then 8 scalar stores:
vextractps $0x1,%xmm0,(%rsi,%rdx,4)     #     AVX2 has no scatter
\end{lstlisting}
The loads are SIMD; the stores cannot be, since AVX2 has no scatter.

\subsection{Interpretation}
The evidence matches the hypothesis: \code{native} uses \code{ymm} for SAXPY, multiply and the
stencil (with FMA in SAXPY and the stencil), \code{autovec} uses \code{xmm}, and \code{scalar} is
genuinely scalar. Tails take a narrower vector step plus a scalar remainder; GCC does not use AVX2's
masked \code{vmaskmovps} (no AVX-512 masks here). The dot product's packed multiplies feed a
latency-bound scalar accumulator, so ``vectorized'' does not mean ``fast'' (\cref{sec:negative}).

% !TeX root = ../main.tex
\section{Scalar vs.\ SIMD Baseline}\label{sec:baseline}
\rubric{Baseline, correctness, reproducibility}{15}{Speedup and achieved throughput for every kernel
across working-set sizes.}

\begin{hypothesis}
  For compute-friendly, unit-stride kernels (SAXPY, elementwise multiply, stencil) we expect the
  \code{native} speedup to approach the float32 lane count ($8\times$) while data is L1/L2-resident,
  and to fall toward $1\times$ once the working set is DRAM-resident and all variants share the same
  bandwidth. \code{autovec} (\num{4} lanes) should track at roughly half the \code{native} speedup.
  The dot product should barely speed up at any size because its reduction does not reassociate
  (\cref{sec:vectorization}).
\end{hypothesis}

\begin{knobs}
  Fixed       & float32; unit stride; aligned arrays; one thread pinned to CPU~2 \\
  Varied      & kernel; variant (scalar, autovec, native); $N = 2^{8} \ldots 2^{25}$ \\
  Repetitions & 12 per point; median, IQR error bars \\
  Script      & \code{SIMD_Advantage_Profiling/scripts/run.sh baseline} \\
\end{knobs}

\begin{figure}[tbp]
  \centering
  \begin{minipage}[t]{0.49\linewidth}
  \IfDataExists{baseline}{%
    \begin{tikzpicture}
      \begin{axis}[acs, xmode=log, log basis x=2, xlabel={$N$ (elements)}, ylabel={Speedup over scalar},
                   width=\linewidth, height=0.68\linewidth, legend style={font=\scriptsize, at={(0.98,0.98)}, anchor=north east}]
        \addplot+[blue, mark=*, mark options={fill=blue}, iqr bars] table[x=n, y=speedup_vs_scalar, speedup iqr, only if={kernel}{saxpy}, only if={variant}{native}] {data/baseline.csv};
        \addlegendentry{SAXPY native}
        \addplot+[blue, densely dashed, mark=o, mark options={solid}, iqr bars] table[x=n, y=speedup_vs_scalar, speedup iqr, only if={kernel}{saxpy}, only if={variant}{autovec}] {data/baseline.csv};
        \addlegendentry{SAXPY autovec}
        \addplot+[red, mark=square*, mark options={fill=red}, iqr bars] table[x=n, y=speedup_vs_scalar, speedup iqr, only if={kernel}{mul}, only if={variant}{native}] {data/baseline.csv};
        \addlegendentry{multiply native}
        \addplot+[brown!80!black, mark=triangle*, mark options={fill=brown!80!black}, iqr bars] table[x=n, y=speedup_vs_scalar, speedup iqr, only if={kernel}{stencil}, only if={variant}{native}] {data/baseline.csv};
        \addlegendentry{stencil native}
        \addplot+[black, mark=diamond*, mark options={fill=black}, iqr bars] table[x=n, y=speedup_vs_scalar, speedup iqr, only if={kernel}{dot}, only if={variant}{native}] {data/baseline.csv};
        \addlegendentry{dot native}
        \acsvline{4096}{L1d}
        \acsvline{32768}{L2}
        \acsvline{2600000}{L3}
      \end{axis}
    \end{tikzpicture}%
  }
  \caption{Speedup of the SIMD builds over scalar (median; error bars: range implied by the two
  IQRs). Gains peak in cache and compress in DRAM; the dot product never benefits.}\label{fig:baseline-speedup}
  \end{minipage}\hfill
  \begin{minipage}[t]{0.49\linewidth}
  \IfDataExists{baseline}{%
    \begin{tikzpicture}
      \begin{axis}[acs, xmode=log, log basis x=2, ymode=log, xlabel={$N$ (elements)}, ylabel={Throughput (elements/s)},
                   width=\linewidth, height=0.68\linewidth, legend pos=north east, legend style={font=\scriptsize}]
        \addplot+[iqr bars] table[x=n, y=elems_per_s_median, iqr of=elems_per_s, only if={kernel}{saxpy}, only if={variant}{scalar}] {data/baseline.csv};
        \addlegendentry{SAXPY scalar}
        \addplot+[iqr bars] table[x=n, y=elems_per_s_median, iqr of=elems_per_s, only if={kernel}{saxpy}, only if={variant}{native}] {data/baseline.csv};
        \addlegendentry{SAXPY native}
        \acsvline{4096}{L1d}
        \acsvline{32768}{L2}
      \end{axis}
    \end{tikzpicture}%
  }
  \caption{Achieved SAXPY throughput (median, IQR). Native drops to a DRAM-limited plateau;
  scalar loses comparatively little outside cache.}\label{fig:baseline-throughput}
  \end{minipage}
\end{figure}

\begin{table}[tbp]
  \centering\small
  \caption{Achieved GFLOP/s (median) of every kernel and build at one $N$ per level: L1
  ($N=\num{1024}$), L2 (\num{16384}), L3 (\num{524288}) and DRAM (\num{16777216}). The same points
  appear in the roofline (\cref{fig:roofline}).}\label{tab:baseline-gflops}
  \IfDataExists{baseline_gflops}{%
    \pgfplotstabletypeset[
      col sep=comma, string type,
      columns={kernel,variant,L1,L2,L3,DRAM},
      columns/kernel/.style={column name=Kernel},
      columns/variant/.style={column name=Build},
      columns/L1/.style={column type=r}, columns/L2/.style={column type=r},
      columns/L3/.style={column type=r}, columns/DRAM/.style={column type=r},
      every head row/.style={before row=\toprule, after row=\midrule},
      every row no 3/.style={before row=\midrule}, every row no 6/.style={before row=\midrule},
      every row no 9/.style={before row=\midrule},
      every last row/.style={after row=\bottomrule},
    ]{data/baseline_gflops.csv}%
  }
\end{table}

\subsection{Analysis}
At $N=\num{4096}$ (L1-resident) the \code{native} speedups are SAXPY $5.79\times$
(\qty{20.2}{\giga\flop\per\second}), elementwise multiply $3.55\times$, stencil $3.49\times$
(\qty{17.3}{\giga\flop\per\second}) and dot $1.01\times$; \code{autovec} reaches $3.02\times$,
$2.54\times$ and $2.79\times$ for the first three. None reaches $8\times$ even in cache: native SAXPY
needs \numrange{0.22}{0.28}~cycles per element in L1 (\cref{sec:locality}) against the
$0.125$~cycles that two load ports and one store port would allow, so load/store issue, not FLOPs,
limits it.

\emph{Elementwise multiply} streams three arrays (\qty{12}{\byte}/element) and leaves L1 earlier
($N\approx\num{2700}$); its speedup peaks at $5.54\times$ at $N=\num{2048}$. The \emph{stencil} does
the most work per element (5~FLOP, reusing $x$) and keeps a $2.6$--$2.7\times$ advantage even in
DRAM. Its cache-resident speedup is irregular: $3.3$--$3.5\times$ for
$N=\num{1024}\ldots\num{16384}$ but $4.5$--$6.2\times$ for $N=\num{32768}\ldots\num{1048576}$
(\cref{sec:anomalies}). The \emph{dot product}, the reduction case of \cref{sec:vectorization},
stays at $\approx1\times$ (\qtyrange{1.33}{1.49}{\giga\flop\per\second}) at every size.

By $N=\num{16777216}$ (\qty{128}{\mebi\byte}) every speedup has compressed
(\cref{fig:baseline-speedup,tab:baseline-gflops}). \Cref{fig:baseline-throughput} shows the
asymmetry: native SAXPY falls from \num{20.2} to \qty{4.7}{\giga\flop\per\second}, scalar only from
about \num{3.5} to \qty{2.4}{\giga\flop\per\second}, because scalar SAXPY stays largely
instruction-bound even in DRAM (\cref{sec:locality}). Other sweeps give $1.88\text{--}2.02\times$ for SAXPY at
$N=2^{24}$--$2^{25}$.

% !TeX root = ../main.tex
\section{Locality Sweep}\label{sec:locality}
\rubric{Locality and memory-bound analysis}{20}{Working-set transitions identified, and the compression
of SIMD gains in the memory-bound regime explained.}

\begin{hypothesis}
  For float32 SAXPY (two arrays, \qty{8}{\byte}/element) the per-element time should be flat and low
  while the working set fits in L1 ($2N\cdot\qty{4}{\byte}\le\qty{32}{\kibi\byte}$, i.e.\
  $N\lesssim\num{4096}$), step up at the L2 boundary ($N\approx\num{32768}$) and again past L3
  ($N\approx\num{2.6e6}$), and saturate in DRAM. The scalar and SIMD curves should be far apart in
  L1 and converge once DRAM bandwidth, which both share, dominates.
\end{hypothesis}

\begin{knobs}
  Fixed       & kernel SAXPY; float32; unit stride; aligned; pinned to CPU~2 \\
  Varied      & $N$ (dense grid, powers of two and midpoints, across L1/L2/LLC/DRAM); variant \\
  Counters    & user-space cycles, instructions, LLC and L1D read misses (\cref{sec:ceilings-method}) \\
  Bandwidth   & elements/s $\times\,\qty{12}{\byte}$ (read $x$, read and write $y$) vs.\ each level's triad roof \\
  Script      & \code{SIMD_Advantage_Profiling/scripts/run.sh locality} \\
\end{knobs}

\begin{figure}[tbp]
  \centering
  \begin{minipage}[t]{0.49\linewidth}
  \IfDataExists{locality}{%
    \begin{tikzpicture}
      \begin{axis}[acs, xmode=log, log basis x=2, ymode=log, xlabel={$N$ (elements)}, ylabel={Time per element (ns)},
                   ytick={0.05,0.1,0.2,0.5,1}, yticklabels={0.05,0.1,0.2,0.5,1}, width=\linewidth, height=0.68\linewidth, legend pos=south east, legend style={font=\scriptsize}]
        \addplot+[iqr bars] table[x=n, y=ns_per_elem_median, iqr of=ns_per_elem, only if={variant}{scalar}] {data/locality.csv};
        \addlegendentry{scalar}
        \addplot+[iqr bars] table[x=n, y=ns_per_elem_median, iqr of=ns_per_elem, only if={variant}{autovec}] {data/locality.csv};
        \addlegendentry{autovec}
        \addplot+[iqr bars] table[x=n, y=ns_per_elem_median, iqr of=ns_per_elem, only if={variant}{native}] {data/locality.csv};
        \addlegendentry{native}
        \acsvline{4096}{L1d}
        \acsvline{32768}{L2}
        \acsvline{2600000}{L3}
      \end{axis}
    \end{tikzpicture}%
  }
  \caption{SAXPY time per element across cache levels: native rises in steps at each
  boundary.}\label{fig:locality}
  \end{minipage}\hfill
  \begin{minipage}[t]{0.49\linewidth}
  \IfDataExists{locality}{%
    \begin{tikzpicture}
      \begin{axis}[acs, xmode=log, log basis x=2, ymode=log, xlabel={Footprint (bytes)},
                   ylabel={Bandwidth (GiB/s)}, ytick={5,10,20,50,100,200,500},
                   yticklabels={5,10,20,50,100,200,500}, width=\linewidth, height=0.68\linewidth, legend pos=south west, legend style={font=\tiny}]
        \addplot+[iqr bars] table[x expr=\thisrow{n}*8, y expr=\thisrow{elems_per_s_median}*12/1073741824,
            y error plus expr=(\thisrow{elems_per_s_q75}-\thisrow{elems_per_s_median})*12/1073741824,
            y error minus expr=(\thisrow{elems_per_s_median}-\thisrow{elems_per_s_q25})*12/1073741824,
            only if={variant}{scalar}] {data/locality.csv};
        \addlegendentry{SAXPY scalar}
        \addplot+[iqr bars] table[x expr=\thisrow{n}*8, y expr=\thisrow{elems_per_s_median}*12/1073741824,
            y error plus expr=(\thisrow{elems_per_s_q75}-\thisrow{elems_per_s_median})*12/1073741824,
            y error minus expr=(\thisrow{elems_per_s_median}-\thisrow{elems_per_s_q25})*12/1073741824,
            only if={variant}{native}] {data/locality.csv};
        \addlegendentry{SAXPY native}
        \IfFileExists{data/ceilings.csv}{%
          \addplot+[no marks, thick, black] table[x=bytes,
              y expr={\thisrow{bytes} < 32768 ? \thisrow{gib_per_s_median} : \thisrow{gib_per_s_total_median}},
              only if={bench}{triad}, only if={label}{sweep}] {data/ceilings.csv};
          \addlegendentry{triad (measured roof)}
          \addplot+[no marks, densely dashed, black!60] table[x=bytes, y=gib_per_s_total_median,
              only if={bench}{read}, only if={label}{sweep}] {data/ceilings.csv};
          \addlegendentry{read (add-latency-bound)}
        }{}
        \acsvline{32768}{L1d}
        \acsvline{262144}{L2}
        \acsvline{20971520}{L3}
      \end{axis}
    \end{tikzpicture}%
  }
  \caption{Effective bandwidth versus footprint: SAXPY against the independently measured triad
  (explicit traffic below \qty{32}{\kibi\byte}, incl.\ write-allocate above). Where native meets the
  triad it is bandwidth-bound. The dashed read stream is add-latency-bound
  ($\approx\qty{89}{\gibi\byte\per\second}$ in L1 and L2), not a ceiling.}\label{fig:locality-bw}
  \end{minipage}
\end{figure}

\begin{table}[tbp]
  \centering\small
  \caption{Effective SAXPY bandwidth (GiB/s, \qty{12}{\byte}/element) at one $N$ per level, next to
  that level's independently measured triad roof (\cref{sec:ceilings-method}).}\label{tab:locality-bw}
  \IfDataExists{locality_levels}{%
    \pgfplotstabletypeset[
      col sep=comma, string type,
      columns={level,n,footprint,scalar,autovec,native,triad},
      columns/level/.style={column name=Level},
      columns/n/.style={column name=$N$},
      columns/footprint/.style={column name=Footprint},
      columns/scalar/.style={column name=scalar, column type=r},
      columns/autovec/.style={column name=autovec, column type=r},
      columns/native/.style={column name=native, column type=r},
      columns/triad/.style={column name=triad roof, column type=r},
      every head row/.style={before row=\toprule, after row=\midrule},
      every last row/.style={after row=\bottomrule},
    ]{data/locality_levels.csv}%
  }
\end{table}



\begin{figure}[tbp]
  \centering
  \IfDataExists{locality_counters}{%
    \begin{subfigure}[t]{0.49\linewidth}
      \begin{tikzpicture}
        \begin{axis}[acs, width=0.95\linewidth, height=0.58\linewidth, xmode=log, log basis x=2, ymode=log,
                     xlabel={$N$}, ylabel={LLC misses / element}, legend pos=north west,
                     legend style={font=\scriptsize}]
          \addplot+[iqr bars] table[x=n, y=llc_miss_per_elem_median, iqr of=llc_miss_per_elem,
              only if={variant}{scalar}] {data/locality_counters.csv};
          \addlegendentry{scalar}
          \addplot+[iqr bars] table[x=n, y=llc_miss_per_elem_median, iqr of=llc_miss_per_elem,
              only if={variant}{native}] {data/locality_counters.csv};
          \addlegendentry{native}
          \acsvline{2600000}{L3}
        \end{axis}
      \end{tikzpicture}
      \caption{LLC misses per element.}
    \end{subfigure}\hfill
    \begin{subfigure}[t]{0.49\linewidth}
      \begin{tikzpicture}
        \begin{axis}[acs, width=0.95\linewidth, height=0.58\linewidth, xmode=log, log basis x=2, ymode=log,
                     xlabel={$N$}, ylabel={Cycles / element}, legend pos=north west,
                     ytick={0.1,0.2,0.5,1,2,5,10}, yticklabels={0.1,0.2,0.5,1,2,5,10},
                     legend style={font=\scriptsize}]
          \addplot+[iqr bars] table[x=n, y=cycles_per_elem_median, iqr of=cycles_per_elem,
              only if={variant}{scalar}] {data/locality_counters.csv};
          \addlegendentry{scalar}
          \addplot+[iqr bars] table[x=n, y=cycles_per_elem_median, iqr of=cycles_per_elem,
              only if={variant}{native}] {data/locality_counters.csv};
          \addlegendentry{native}
          \acsvline{4096}{L1d}
          \acsvline{32768}{L2}
          \acsvline{2600000}{L3}
        \end{axis}
      \end{tikzpicture}
      \caption{Core cycles per element.}
    \end{subfigure}%
  }
  \caption{User-space hardware counters for the locality sweep (\code{perf_event_open}).}\label{fig:locality-counters}
\end{figure}

\subsection{Analysis}
The transitions line up with the cache sizes. In L1 native SAXPY takes
\qtyrange{0.080}{0.100}{\nano\second} per element, $5.6\text{--}6.9\times$ faster than scalar; the
speedup drops to $4.0\text{--}4.2\times$ past L1, $3.1\times$ past L2 and $2.8\times$ on the L3
plateau. Past L3 the gap narrows but does \emph{not} close, contradicting the hypothesis: native
takes \qtyrange{0.43}{0.45}{\nano\second} and scalar \qtyrange{0.83}{0.85}{\nano\second}, a speedup
of $1.88\text{--}1.96\times$.

In bandwidth terms (\cref{tab:locality-bw,fig:locality-bw}) native SAXPY reaches 76, 86, 100 and
\SI{103}{\percent} of the measured triad roof in L1, L2, L3 and DRAM: bandwidth-bound in L3 and DRAM
(in DRAM still below the \qty{25.7}{\gibi\byte\per\second} of the copy stream, which has the same
two-reads-one-write pattern) and issue-bound in L1 and L2. Scalar is flat at
\qtyrange{17}{20.5}{\gibi\byte\per\second} through L3, far below every roof.

\emph{Counters} (\cref{fig:locality-counters}). For DRAM-resident native SAXPY, LLC misses reach
\numrange{0.122}{0.128} per element, the compulsory $2/16 = 0.125$ (one line each of $x$ and $y$ per
\num{16} elements), so the vector loop moves no more data than it must. Scalar shows
\numrange{0.156}{0.164}, about \SI{28}{\percent} more, for a reason we have not established
(\cref{sec:anomalies}). The cycle counts explain
the remaining gap:

\begin{center}\small
  \begin{tabular}{@{}lrrrr@{}}
    \toprule
    cycles/element (IPC) & L1 ($N=\num{2048}$) & L2 ($N=\num{16384}$) & L3 ($N=\num{524288}$) & DRAM ($N=\num{33554432}$) \\
    \midrule
    scalar & 1.53 (4.6) & 1.68 (4.2) & 1.81 (3.9) & 2.34 (3.0) \\
    native & 0.26 (3.1) & 0.40 (1.9) & 0.64 (1.2) & 1.25 (0.6) \\
    \bottomrule
  \end{tabular}
\end{center}
Native SAXPY's cycles per element grow $4.7\times$ from L1 to DRAM and its IPC falls to \num{0.6}:
it waits on memory. Scalar already spends \num{1.53} cycles per element in L1 at IPC \num{4.6}, this
core's issue limit, and in DRAM still retires \num{3.0} instructions per cycle: most of its time is
execution. $2.34/1.25 = 1.88$ is exactly the measured DRAM speedup at this size. So once
native sits on the bandwidth roof, the speedup is scalar's instruction-bound time over the memory
time, a ratio that depends on the clock (\cref{sec:method}).

% !TeX root = ../main.tex
\section{Alignment and Tail Handling}\label{sec:alignment}
\rubric{Alignment, tail, stride/gather studies}{20}{Quantified alignment, prologue/epilogue and masking
effects with a convincing explanation (with \cref{sec:stride}).}

\begin{hypothesis}
  A start address that is not \qty{32}{\byte}-aligned should make some \qty{256}{bit} loads
  unaligned and/or trigger a scalar peel prologue, a modest fixed cost most visible at small $N$.
  On a client Skylake core unaligned \code{vmovups} within one cache line are nearly free, so we
  expect a small penalty (tens of percent) that recurs with \qty{64}{\byte} period across a line.
\end{hypothesis}
\begin{hypothesis}
  (For the extended sweep.) \code{scalar}'s \qty{4}{\byte} loads never straddle a line, so it should
  be insensitive to the offset, and \code{autovec} (\qty{16}{\byte} loads) should suffer less than
  \code{native}. In DRAM the misalignment cost should vanish into memory stalls. For short loops the
  time per \emph{call} should grow in steps of \num{8} elements (\code{native}) or \num{4}
  (\code{autovec}), with the vectorized epilogue making the remainder cheap but not free.
\end{hypothesis}

\begin{knobs}
  Fixed       & kernel SAXPY; float32; unit stride \\
  Varied      & start offsets $0, 4, 8, \ldots, \qty{60}{\byte}$ (every multiple of \num{4}) at $N=\num{4096}$
                (L1-resident, isolates the prologue from bandwidth) and $N=\num{16777216}$
                (\qty{128}{\mebi\byte}, DRAM-resident); tail: every $N$ from \num{8} to \num{72} and
                $N$ stepping across an \num{8}-lane boundary near \num{4096}, offset $0$; all three
                build variants \\
  Script      & \code{SIMD_Advantage_Profiling/scripts/run.sh alignment} \\
\end{knobs}

\begin{figure}[tbp]
  \centering
  \begin{subfigure}[t]{0.49\linewidth}
    \begin{tikzpicture}
      \begin{axis}[acs, height=0.56\linewidth, xlabel={Start offset (bytes)}, ylabel={Time per element (ns)},
                   legend style={font=\scriptsize}, legend pos=north east, ymin=0]
        \acsvariantlegend
        \addplot+[var scalar, forget plot, iqr bars] table[x=offset, y=ns_per_elem_median, iqr of=ns_per_elem, only if={variant}{scalar}, only if={n}{4096}] {data/alignment.csv};
        \addplot+[var autovec, forget plot, iqr bars] table[x=offset, y=ns_per_elem_median, iqr of=ns_per_elem, only if={variant}{autovec}, only if={n}{4096}] {data/alignment.csv};
        \addplot+[var native, forget plot, iqr bars] table[x=offset, y=ns_per_elem_median, iqr of=ns_per_elem, only if={variant}{native}, only if={n}{4096}] {data/alignment.csv};
      \end{axis}
    \end{tikzpicture}
    \caption{$N=\num{4096}$ (L1-resident).}
  \end{subfigure}\hfill
  \begin{subfigure}[t]{0.49\linewidth}
    \begin{tikzpicture}
      \begin{axis}[acs, height=0.56\linewidth, xlabel={Start offset (bytes)}, ylabel={Time per element (ns)}, ymin=0]
        \addplot+[var scalar, forget plot, iqr bars] table[x=offset, y=ns_per_elem_median, iqr of=ns_per_elem, only if={variant}{scalar}, only if={n}{16777216}] {data/alignment.csv};
        \addplot+[var autovec, forget plot, iqr bars] table[x=offset, y=ns_per_elem_median, iqr of=ns_per_elem, only if={variant}{autovec}, only if={n}{16777216}] {data/alignment.csv};
        \addplot+[var native, forget plot, iqr bars] table[x=offset, y=ns_per_elem_median, iqr of=ns_per_elem, only if={variant}{native}, only if={n}{16777216}] {data/alignment.csv};
      \end{axis}
    \end{tikzpicture}
    \caption{$N=\num{16777216}$ (DRAM-resident).}
  \end{subfigure}
  \caption{Cost of a misaligned start for SAXPY, by build variant (median, IQR error bars).}\label{fig:alignment}
\end{figure}

\begin{figure}[tbp]
  \centering
  \begin{subfigure}[t]{0.49\linewidth}
    \begin{tikzpicture}
      \begin{axis}[acs, height=0.56\linewidth, xlabel={$N$ (elements)}, ylabel={Time per call (ns)},
                   restrict x to domain=0:100, xmin=6, xmax=74, xtick={8,16,...,72}, legend style={font=\scriptsize}, legend pos=north west]
        \acsvariantlegend
        \addplot+[mark size=1pt, var scalar, forget plot, iqr bars] table[x=n, y expr=\thisrow{ns_per_elem_median}*\thisrow{n}, y error plus expr=(\thisrow{ns_per_elem_q75}-\thisrow{ns_per_elem_median})*\thisrow{n}, y error minus expr=(\thisrow{ns_per_elem_median}-\thisrow{ns_per_elem_q25})*\thisrow{n}, only if={variant}{scalar}] {data/alignment_tail.csv};
        \addplot+[mark size=1pt, var autovec, forget plot, iqr bars] table[x=n, y expr=\thisrow{ns_per_elem_median}*\thisrow{n}, y error plus expr=(\thisrow{ns_per_elem_q75}-\thisrow{ns_per_elem_median})*\thisrow{n}, y error minus expr=(\thisrow{ns_per_elem_median}-\thisrow{ns_per_elem_q25})*\thisrow{n}, only if={variant}{autovec}] {data/alignment_tail.csv};
        \addplot+[mark size=1pt, var native, forget plot, iqr bars] table[x=n, y expr=\thisrow{ns_per_elem_median}*\thisrow{n}, y error plus expr=(\thisrow{ns_per_elem_q75}-\thisrow{ns_per_elem_median})*\thisrow{n}, y error minus expr=(\thisrow{ns_per_elem_median}-\thisrow{ns_per_elem_q25})*\thisrow{n}, only if={variant}{native}] {data/alignment_tail.csv};
      \end{axis}
    \end{tikzpicture}
    \caption{Short loops: time per call, $N=8\ldots72$ (IQR bars, mostly hidden by the marks).}
  \end{subfigure}\hfill
  \begin{subfigure}[t]{0.49\linewidth}
    \begin{tikzpicture}
      \begin{axis}[acs, height=0.56\linewidth, xlabel={$N$ (elements)}, ylabel={Time per element (ns)},
                   restrict x to domain=4080:4140, xmin=4084, xmax=4132, ymin=0]
        \addplot+[var scalar, forget plot, iqr bars] table[x=n, y=ns_per_elem_median, iqr of=ns_per_elem, only if={variant}{scalar}] {data/alignment_tail.csv};
        \addplot+[var autovec, forget plot, iqr bars] table[x=n, y=ns_per_elem_median, iqr of=ns_per_elem, only if={variant}{autovec}] {data/alignment_tail.csv};
        \addplot+[var native, forget plot, iqr bars] table[x=n, y=ns_per_elem_median, iqr of=ns_per_elem, only if={variant}{native}] {data/alignment_tail.csv};
      \end{axis}
    \end{tikzpicture}
    \caption{Across an \num{8}-lane boundary near $N=\num{4096}$.}
  \end{subfigure}
  \caption{Tail handling: remainder elements processed by the vectorized epilogue (offset $0$).}\label{fig:tail}
\end{figure}


\subsection{Analysis}
\emph{Misaligned starts, L1-resident} (\cref{fig:alignment}a). Native SAXPY runs at
\qtyrange{0.098}{0.100}{\nano\second} per element at the \qty{32}{\byte}-aligned offsets $0$ and
\num{32}, \SI{46}{\percent} slower at most other offsets and up to \SI{60}{\percent} slower at
\numlist{28;36;40;60}. \code{autovec} is fast exactly at multiples of \num{16} and
\SIrange{13}{25}{\percent} slower elsewhere; \code{scalar} varies by only \SI{1.0}{\percent}. This is
what cache-line splits predict: off a 32-byte boundary every second \qty{32}{\byte} load straddles a
\qty{64}{\byte} line, every fourth \qty{16}{\byte} load does, and a \qty{4}{\byte} load never does;
GCC does not peel to an aligned start, so every split pays. The second level at
\numlist{28;36;40;60} is unexplained (untested: split stores to \code{y} interacting with loads of
\code{x}). Every IQR is below \SI{0.3}{\percent}.

\emph{DRAM-resident} (\cref{fig:alignment}b). The spread across offsets is only \SI{6.1}{\percent}
(native), \SI{5.7}{\percent} (autovec) and \SI{2.1}{\percent} (scalar), with no split-load pattern:
the cost hides under DRAM misses, as hypothesized.

\emph{Tails} (\cref{fig:tail}). Near $N=\num{4096}$ the \numrange{1}{7} leftover elements are
invisible. In short loops ($N=8\ldots72$) native's time per call is lowest at multiples of \num{8}
(\qty{13.7}{\nano\second} at $N=16$ to \qty{16.7}{\nano\second} at $N=72$) and typically
\qtyrange{2}{3}{\nano\second} higher at $N\equiv7 \pmod 8$: the sawtooth of one \qty{128}{bit} step
plus up to three scalar iterations (\cref{sec:vectorization}). A fixed $\approx\qtyrange{10}{13}{\nano\second}$ per call (call, setup, aliasing check)
means native beats scalar by only $1.11\times$ at $N=8$ and $3.32\times$ at $N=72$. A masked AVX2
tail (\code{vmaskmovps}) would replace the remainder path; we did not implement one.

% !TeX root = ../main.tex
\section{Stride and Gather Access}\label{sec:stride}

\begin{hypothesis}
  Unit stride lets each \qty{64}{\byte} line serve \num{16} float32 elements; at stride \num{16} each
  element needs its own line. Time per processed element should rise roughly linearly with stride
  until it saturates at one line per element, and the SIMD advantage should vanish for any
  non-unit stride, because lanes must be loaded individually. Gather should be the worst case.
\end{hypothesis}
\begin{hypothesis}
  (For the extended sweep.) The explicit AVX2 gather (\code{vgatherdps}, \cref{sec:vectorization})
  still needs one line per element and eight scalar stores, so it should be at best modestly faster
  than the scalar gather in cache. This CPU runs the \enquote{Gather Data Sampling} microcode
  mitigation (\code{/sys/devices/system/cpu/vulnerabilities/gather_data_sampling}), reported to slow
  \code{vgatherdps} considerably, so it may even lose.
\end{hypothesis}

\begin{knobs}
  Fixed       & SAXPY; float32 \\
  Varied      & access (unit, stride $1/2/4/8/16$, gather, explicit AVX2 gather); variant (scalar,
                autovec, native; AVX2 gather native only); $N=\num{32768}$ processed elements
                (span $\qty{256}{\kibi\byte}\times s$: L2 for $s=1$, L3 for $s=16$) and
                $N=\num{4194304}$ (DRAM for every $s$) \\
  Definition  & stride $s$ processes elements $0, s, 2s, \ldots$ over an array of $N s$ elements;
                gather applies SAXPY through a random permutation index so every element is touched
                once with no spatial locality. \\
  Script      & \code{SIMD_Advantage_Profiling/scripts/run.sh stride} \\
\end{knobs}

\begin{figure}[tbp]
  \centering
  \IfDataExists{stride}{%
    \begin{subfigure}[t]{0.49\linewidth}
      \begin{tikzpicture}
        \begin{axis}[acs, height=0.56\linewidth, xmode=log, log basis x=2, ymode=log, xlabel={Stride (elements)},
                     ylabel={Time per processed element (ns)}, legend pos=north west, legend style={font=\scriptsize},
                     ytick={0.1,0.2,0.5,1,2,5,10}, yticklabels={0.1,0.2,0.5,1,2,5,10}]
          \acsvariantlegend
          \addplot+[var scalar, forget plot, iqr bars] table[x=stride, y=ns_per_elem_median, iqr of=ns_per_elem, only if={access}{stride}, only if={variant}{scalar}, only if={n}{32768}] {data/stride.csv};
          \addplot+[var autovec, forget plot, iqr bars] table[x=stride, y=ns_per_elem_median, iqr of=ns_per_elem, only if={access}{stride}, only if={variant}{autovec}, only if={n}{32768}] {data/stride.csv};
          \addplot+[var native, forget plot, iqr bars] table[x=stride, y=ns_per_elem_median, iqr of=ns_per_elem, only if={access}{stride}, only if={variant}{native}, only if={n}{32768}] {data/stride.csv};
        \end{axis}
      \end{tikzpicture}
      \caption{$N=\num{32768}$ (cache-resident span).}
    \end{subfigure}\hfill
    \begin{subfigure}[t]{0.49\linewidth}
      \begin{tikzpicture}
        \begin{axis}[acs, height=0.56\linewidth, xmode=log, log basis x=2, ymode=log, xlabel={Stride (elements)},
                     ylabel={Time per processed element (ns)},
                     ytick={0.1,0.2,0.5,1,2,5,10}, yticklabels={0.1,0.2,0.5,1,2,5,10}]
          \addplot+[var scalar, forget plot, iqr bars] table[x=stride, y=ns_per_elem_median, iqr of=ns_per_elem, only if={access}{stride}, only if={variant}{scalar}, only if={n}{4194304}] {data/stride.csv};
          \addplot+[var autovec, forget plot, iqr bars] table[x=stride, y=ns_per_elem_median, iqr of=ns_per_elem, only if={access}{stride}, only if={variant}{autovec}, only if={n}{4194304}] {data/stride.csv};
          \addplot+[var native, forget plot, iqr bars] table[x=stride, y=ns_per_elem_median, iqr of=ns_per_elem, only if={access}{stride}, only if={variant}{native}, only if={n}{4194304}] {data/stride.csv};
        \end{axis}
      \end{tikzpicture}
      \caption{$N=\num{4194304}$ (DRAM-resident).}
    \end{subfigure}%
  }
  \caption{Effect of stride on SAXPY time per \emph{processed} element (median, IQR error bars). The
  three builds coincide for every non-unit stride.}\label{fig:stride}
\end{figure}

\begin{table}[tbp]
  \centering\small
  \caption{Time per processed element (ns) by access pattern. \emph{AVX2 gather} is the explicit
  \code{vgatherdps} kernel (native build only); speedup is scalar over the faster native
  kernel (the AVX2 gather has no strided form, hence --).}\label{tab:stride}
  \IfDataExists{stride_summary}{%
    \pgfplotstabletypeset[
      col sep=comma, string type,
      columns={n,pattern,scalar,autovec,native,native_avx2,speedup},
      columns/n/.style={column name=$N$},
      columns/pattern/.style={column name=Access},
      columns/scalar/.style={column name=scalar, column type=r},
      columns/autovec/.style={column name=autovec, column type=r},
      columns/native/.style={column name=native, column type=r},
      columns/native_avx2/.style={column name=AVX2 gather, column type=r},
      columns/speedup/.style={column name=speedup ($\times$), column type=r},
      every head row/.style={before row=\toprule, after row=\midrule},
      every row no 6/.style={before row=\midrule},
      every last row/.style={after row=\bottomrule},
    ]{data/stride_summary.csv}%
  }
\end{table}

\subsection{Analysis}
Only unit stride lets \code{native} beat \code{scalar} (\cref{tab:stride}): $3.6\text{--}4.1\times$
at $N=\num{32768}$ and $2.35\times$ in DRAM. The range comes from two runs of identical machine
code in different processes (the table's \emph{unit} row and the stride-1 point of
\cref{fig:stride}a; \cref{sec:anomalies}).

At stride~2 and beyond the builds agree within \numrange{0.98}{1.03}$\times$, because every $s>1$
runs the scalar copy of the versioned loop (\cref{sec:vectorization}), and cost follows the cache
lines touched. In DRAM native time per element grows \num{1.13} $\to$ \num{1.85} $\to$ \num{3.61}
$\to$ \qty{7.34}{\nano\second} from stride~2 to~16. Counting fetched traffic as
$3\times\min(4s,64)$~\unit{\byte} per element (\code{stride_bw.csv}), useful bandwidth falls from
\num{9.9} to \qty{1.5}{\gibi\byte\per\second} while fetched traffic stays at
\numrange{24.1}{24.8}~\unit{\gibi\byte\per\second} from stride~4 on, the DRAM triad roof: memory is
saturated with lines of which only $4s/64$ is used. Gather is the worst pattern:
\qty{10.2}{\nano\second} per element in DRAM, $31\times$ unit stride, because random order also
defeats the prefetcher.

\emph{Hardware gather.} The explicit AVX2 kernel takes $1.59\times$ (\qty{3.23}{\nano\second},
$N=\num{32768}$) and $1.07\times$ (\qty{10.91}{\nano\second}, DRAM) the time of the scalar build's
gather, as the second hypothesis anticipated: two mitigated \code{vgatherdps} plus eight scalar
stores per step. Testing without the mitigation needs a reboot as root. Access pattern, not vector
width, is the first-order knob for these kernels.

% !TeX root = ../main.tex
\section{Data Type Comparison}\label{sec:dtype}

\begin{hypothesis}
  A \qty{256}{bit} register holds \num{8} float32 but only \num{4} float64 lanes, so float64 should
  show roughly half the SIMD speedup of float32 for the same kernel while data is cache-resident.
  int32 elementwise multiply should behave like float32 (\num{8} lanes, one \code{vpmulld}). In DRAM,
  float64 moves twice the bytes per element, so its \emph{absolute} throughput should be lower and
  its speedup should compress toward $1\times$ just as float32 does.
\end{hypothesis}

\begin{knobs}
  Fixed       & unit stride; aligned \\
  Varied      & data type (float32, float64, int32); kernel; variant; $N = 2^{10}\ldots2^{24}$ \\
  Script      & \code{SIMD_Advantage_Profiling/scripts/run.sh dtype} \\
\end{knobs}

\begin{figure}[tbp]
  \centering
  \begin{minipage}[t]{0.42\linewidth}
  \IfDataExists{dtype}{%
    \begin{tikzpicture}
      \begin{axis}[acs, xmode=log, log basis x=2, xlabel={$N$ (elements)}, ylabel={Speedup over scalar}, width=\linewidth, height=0.72\linewidth, legend pos=north east, legend style={font=\scriptsize}]
        \addplot+[iqr bars] table[x=n, y=speedup_vs_scalar, speedup iqr, only if={kernel}{mul}, only if={variant}{native}, only if={dtype}{f32}] {data/dtype.csv};
        \addlegendentry{multiply float32}
        \addplot+[iqr bars] table[x=n, y=speedup_vs_scalar, speedup iqr, only if={kernel}{mul}, only if={variant}{native}, only if={dtype}{f64}] {data/dtype.csv};
        \addlegendentry{multiply float64}
        \addplot+[iqr bars] table[x=n, y=speedup_vs_scalar, speedup iqr, only if={kernel}{mul}, only if={variant}{native}, only if={dtype}{i32}] {data/dtype.csv};
        \addlegendentry{multiply int32}
      \end{axis}
    \end{tikzpicture}%
  }
  \caption{Native speedup of elementwise multiply by data type (median; error bars: range implied
  by the two IQRs). float32 and int32 (\num{8} lanes) track each other; float64 (\num{4} lanes) is
  about half. SAXPY and stencil: \cref{tab:dtype}.}\label{fig:dtype}
  \end{minipage}\hfill
  \begin{minipage}[t]{0.56\linewidth}
  \centering\footnotesize\setlength{\tabcolsep}{2.5pt}
  \captionof{table}{Speedup over scalar by kernel and data type (\code{dtype.csv}), at $N=\num{1024}$ (L1),
  \num{16384} (L2) and \num{16777216} (DRAM).}\label{tab:dtype}
  \IfDataExists{dtype_summary}{%
    \pgfplotstabletypeset[
      col sep=comma, string type,
      columns={kernel,dtype,autovec_1k,native_1k,native_16k,native_16m},
      columns/kernel/.style={column name=Kernel},
      columns/dtype/.style={column name=Type},
      columns/autovec_1k/.style={column name={autovec, \num{1024}}, column type=r},
      columns/native_1k/.style={column name={native, \num{1024}}, column type=r},
      columns/native_16k/.style={column name={\num{16384}}, column type=r},
      columns/native_16m/.style={column name={\num{16777216}}, column type=r},
      every head row/.style={before row=\toprule, after row=\midrule},
      every row no 2/.style={before row=\midrule}, every row no 5/.style={before row=\midrule},
      every last row/.style={after row=\bottomrule},
    ]{data/dtype_summary.csv}%
  }
  \end{minipage}
\end{figure}



\subsection{Analysis}
At $N=\num{1024}$ native float64 keeps \numrange{0.52}{0.59} of float32's speedup for all three
kernels (multiply $2.82\times$ vs $5.19\times$, SAXPY $3.28\times$ vs $5.60\times$, stencil
$1.73\times$ vs $3.35\times$), close to the predicted $1/2$, and int32 multiply (one \code{vpmulld}
over \num{8} lanes) lands on the float32 curve ($5.37\times$). The ratio is not universal: at
$N=\num{16384}$ the stencil keeps $3.02\times$ in float64 against $3.44\times$ in float32, because
the float32 stencil is depressed at these sizes (\cref{sec:anomalies}), and by $N=\num{65536}$
multiply narrows to $1.40\times$ vs $2.11\times$.

\emph{Operation cost.} \code{autovec} gains $3.53\times$ on float32 multiply at $N=\num{1024}$ but
only $1.86\times$ on int32: SSE2 has no packed 32-bit multiply (\code{pmulld} is SSE4.1), so GCC
emulates it with seven instructions per \num{4} lanes (two \code{pmuludq}, two shifts, two
\code{pshufd} and a \code{punpckldq}).

In DRAM ($N=\num{16777216}$) every speedup compresses as predicted (native float32
$1.60$--$2.79\times$, float64 $1.17$--$1.68\times$), and float64 keeps \SIrange{60}{73}{\percent} of
the float32 gain. Float32 multiply's $1.51\times$ in the baseline sweep differs by a
\SI{6}{\percent} run-to-run spread in native throughput, so we quote $1.5\text{--}1.6\times$. Float64's absolute multiply
throughput is under half of float32's (\num{0.80} vs \qty{1.72}{\giga\flop\per\second}) because
each element moves twice the bytes. Lane count and operation cost set the cache-resident speedup;
bytes per element set the DRAM throughput.

% !TeX root = ../main.tex
\section{Roofline Interpretation}\label{sec:roofline}
\rubric{Roofline and quantitative insight}{20}{Correct arithmetic intensity, use of measured memory
bandwidth, and a justified compute- vs.\ memory-bound conclusion.}

\begin{hypothesis}
  All four kernels have very low arithmetic intensity (well under \qty{1}{\flop\per\byte}). With
  single-core bandwidth of tens of \unit{\giga\byte\per\second} and a peak vector rate near
  \qty{100}{\giga\flop\per\second}, all should sit on the \emph{bandwidth} slope when DRAM-resident.
  The stencil, with the highest intensity, should come closest to the compute roof.
\end{hypothesis}
\begin{hypothesis}
  (For the measured ceilings.) Even the L1 roof should put the ridge point above
  \qty{0.5}{\flop\per\byte}, so cache-resident native kernels should also be bounded by bandwidth
  (or the load/store ports behind it), not by the FMA units.
\end{hypothesis}

Attainable performance at arithmetic intensity $I$ (FLOPs per byte of memory traffic) is bounded by
\begin{equation}
  P(I) = \min\bigl(P_{\text{peak}},\; B_{\ell} \cdot I\bigr),
\end{equation}
where $B_{\ell}$ is the bandwidth of the memory level $\ell$ holding the data
\cite{williams2009roofline}. We draw one roof per level (a cache-aware roofline); both ceilings
come from separate microbenchmarks (\cref{sec:ceilings-method}), never from the kernels plotted.

\begin{table}[tbp]
  \centering\small
  \caption{Measured roofline ceilings, single core (\code{run.sh ceilings}). Top: register-only
  accumulator loops; FLOP/cycle and clock from the cycle counter. Bottom: STREAM-style bandwidth
  (GiB/s) at one footprint per level; \enquote{incl.\ WA} adds the write-allocate read. Roofs use
  the explicit triad in L1 and triad incl.\ WA below it; the read stream ($^{*}$) is
  add-latency-bound and not a roof.}\label{tab:ceilings}
  \IfDataExists{ceilings_peak}{%
    \pgfplotstabletypeset[
      col sep=comma,
      columns={isa,gflops_median,flops_per_cycle,ghz},
      columns/isa/.style={column name=Peak loop, string type, string replace*={_}{\_}},
      columns/gflops_median/.style={column name=GFLOP/s, fixed, zerofill, precision=1},
      columns/flops_per_cycle/.style={column name=FLOP/cycle, fixed, zerofill, precision=2},
      columns/ghz/.style={column name=GHz, fixed, zerofill, precision=2},
      every head row/.style={before row=\toprule, after row=\midrule},
      every last row/.style={after row=\bottomrule},
    ]{data/ceilings_peak.csv}%
  }\par\medskip
  \IfDataExists{ceilings_bw}{%
    \pgfplotstabletypeset[
      col sep=comma,
      columns={level,bytes,read,copy,triad,triad_total},
      columns/level/.style={column name=Level, string type},
      columns/bytes/.style={column name=Footprint (B), int detect},
      columns/read/.style={column name=read$^{*}$, fixed, zerofill, precision=1},
      columns/copy/.style={column name=copy, fixed, zerofill, precision=1},
      columns/triad/.style={column name=triad, fixed, zerofill, precision=1},
      columns/triad_total/.style={column name={triad incl.\ WA}, fixed, zerofill, precision=1},
      every head row/.style={before row={\toprule & & \multicolumn{4}{c}{GiB/s}\\}, after row=\midrule},
      every last row/.style={after row=\bottomrule},
    ]{data/ceilings_bw.csv}%
  }
\end{table}

\begin{table}[tbp]
  \centering\small
  \caption{Arithmetic intensity (float32) from the actual traffic per element, the same convention
  as the triad roof: L2, L3 and DRAM include the write-allocate read of a store to an unread line;
  in L1 every store hits, so there is none.}\label{tab:roofline-inputs}
  \begin{tabularx}{\linewidth}{@{}lrrX@{}}
    \toprule
    Kernel & $I$, L2--DRAM & $I$, L1 & Traffic per element (L2--DRAM) \\
    \midrule
    SAXPY        & \num{0.167} & \num{0.167} & $2\,\text{FLOP} / \qty{12}{\byte}$: read $x,y$, write $y$ (already read: no write-allocate) \\
    multiply     & \num{0.063} & \num{0.083} & $1\,\text{FLOP} / \qty{16}{\byte}$: read $x,y$, write $z$ + write-allocate of $z$ \\
    stencil      & \num{0.417} & \num{0.625} & $5\,\text{FLOP} / \qty{12}{\byte}$: read $x$ once (reuse), write $y$ + write-allocate \\
    dot product  & \num{0.25}  & \num{0.25}  & $2\,\text{FLOP} / \qty{8}{\byte}$: read $x,y$, no store \\
    triad (roof) & \num{0.125} & \num{0.167} & $2\,\text{FLOP} / \qty{16}{\byte}$: read $b,c$, write $a$ + write-allocate \\
    \bottomrule
  \end{tabularx}
\end{table}

\begin{figure}[tbp]
  \centering
  \IfDataExists{roofline_points}{%
    \begin{tikzpicture}
      \begin{axis}[acs, xmode=log, ymode=log, xlabel={Arithmetic intensity (FLOP/byte)},
                   ylabel={Performance (GFLOP/s)}, width=0.68\linewidth, height=0.42\linewidth,
                   xmin=0.02, xmax=20, ymin=0.5, ymax=400,
                   legend pos=outer north east, legend style={font=\scriptsize}]
        \IfFileExists{data/roofline_roofs.csv}{%
          \addplot[thick, black] table[x=ai, y=gflops, only if={roof}{L1}] {data/roofline_roofs.csv};
          \addlegendentry{L1 roof}
          \addplot[thick, black!70, densely dashed] table[x=ai, y=gflops, only if={roof}{L2}] {data/roofline_roofs.csv};
          \addlegendentry{L2 roof}
          \addplot[thick, black!55, dashdotted] table[x=ai, y=gflops, only if={roof}{L3}] {data/roofline_roofs.csv};
          \addlegendentry{L3 roof}
          \addplot[thick, black!40, dotted] table[x=ai, y=gflops, only if={roof}{DRAM}] {data/roofline_roofs.csv};
          \addlegendentry{DRAM roof}
          \addplot[thin, red!60] table[x=ai, y=gflops, only if={roof}{peak_sse_f32}] {data/roofline_roofs.csv};
          \addlegendentry{SSE peak (autovec)}
          \addplot[thin, red!60, densely dashed] table[x=ai, y=gflops, only if={roof}{peak_scalar_f32}] {data/roofline_roofs.csv};
          \addlegendentry{scalar peak}
        }{}
        \addplot+[only marks, mark=*, nodes near coords, point meta=explicit symbolic,
                  every node near coord/.style={font=\tiny, anchor=south west}, iqr bars]
          table[x=ai, y=gflops_median, iqr of=gflops, meta=kernel, only if={variant}{native}, only if={level}{L1}] {data/roofline_points.csv};
        \addlegendentry{native, L1-resident}
        \addplot+[only marks, mark=square*, iqr bars] table[x=ai, y=gflops_median, iqr of=gflops, only if={variant}{native}, only if={level}{L2}] {data/roofline_points.csv};
        \addlegendentry{native, L2-resident}
        \addplot+[only marks, mark=triangle*, iqr bars] table[x=ai, y=gflops_median, iqr of=gflops, only if={variant}{native}, only if={level}{L3}] {data/roofline_points.csv};
        \addlegendentry{native, L3-resident}
        \addplot+[only marks, mark=diamond*, iqr bars] table[x=ai, y=gflops_median, iqr of=gflops, only if={variant}{native}, only if={level}{DRAM}] {data/roofline_points.csv};
        \addlegendentry{native, DRAM}
        \addplot+[only marks, mark=o, black, iqr bars] table[x=ai, y=gflops_median, iqr of=gflops, only if={variant}{scalar}, only if={level}{DRAM}] {data/roofline_points.csv};
        \addlegendentry{scalar, DRAM}
      \end{axis}
    \end{tikzpicture}%
  }
  \caption{Cache-aware roofline for one core. Roofs: measured triad bandwidth per level, capped
  by the measured AVX2-FMA peak; thin red lines: measured SSE and scalar peaks. Points: native (and
  DRAM scalar) GFLOP/s from \cref{tab:baseline-gflops}, L1 points at their L1 intensity; IQR bars
  mostly hidden by the markers.}\label{fig:roofline}
\end{figure}

\subsection{Analysis}
\emph{Ceilings} (\cref{tab:ceilings}). The AVX2 FMA loop reaches \qty{89.2}{\giga\flop\per\second}
float32 (\num{44.6} float64), \num{31.99}~FLOP/cycle at \qty{2.79}{\giga\hertz}: exactly the
$2~\text{FMA ports} \times 8~\text{lanes} \times 2 = 32$ the core allows. The SSE and scalar loops
also reach their port limits (\num{8.00} and \num{1.99}~FLOP/cycle). The ridge points $P_{\text{peak}}/B_\ell$ are $\approx\num{0.53}$, \num{0.92},
\num{1.71} and \qty{3.42}{\flop\per\byte} from L1 to DRAM; the L1 ridge is just above the
\qty{0.5}{\flop\per\byte} of the second hypothesis, and only the L1 stencil ($I=\num{0.625}$) lies
right of it.

\emph{DRAM-resident kernels.} Every intensity (\numrange{0.06}{0.42}~FLOP/byte) is far left of the
DRAM ridge, so all four kernels are bandwidth-bound, as hypothesized. Native multiply reaches
\SI{100}{\percent} of the DRAM roof, the stencil \SI{101}{\percent} and SAXPY \SI{107}{\percent}
(\num{4.65} vs \qty{4.35}{\giga\flop\per\second}). Sustainable bandwidth depends a little on the
access pattern (SAXPY writes a line it has just read), so the triad roof is a reference, not a bound
to the byte; since it comes from a different loop, landing on it is evidence, not a tautology. The scalar points sit at \SIrange{37}{66}{\percent} of the
roof (dot: \SI{20}{\percent}): the instruction-bound loops of \cref{sec:locality}.

\emph{Cache-resident kernels.} L3-resident native points sit at \SIrange{95}{100}{\percent} of the
L3 roof. In L2 SAXPY and multiply reach \SI{86}{\percent} and \SI{88}{\percent} (stencil
\SI{42}{\percent}); in L1 SAXPY reaches \SI{72}{\percent} of the roof and \SI{23}{\percent} of the FMA
peak, and the stencil \SI{19}{\percent} of the peak that bounds it. So, confirming the second hypothesis,
cache-resident kernels are bounded by load/store and instruction issue, not the FMA units. From L2 outward each step in
\cref{fig:locality} is a move to a lower roof.

\emph{The dot product} is far \emph{below} every roof because it is latency-bound: one scalar
\code{vaddss} per element with a \num{4}-cycle latency allows $0.5$~FLOP/cycle,
\qty{1.39}{\giga\flop\per\second}, and it measures \qtyrange{1.39}{1.42}{\giga\flop\per\second} for
$N=\num{1024}$ to \num{1048576}. Wider vectors cannot help a kernel waiting on memory; only
higher intensity (fusing passes, cache blocking) can.

% !TeX root = ../main.tex
\section{Negative Example}\label{sec:negative}
\rubric{AI-use transparency and failure case}{10}{A meaningful negative or misleading example, verified
and explained (with \cref{sec:ai-use}).}

\begin{hypothesis}
  A naive benchmark---or a quickly generated one that only ever tests a small array---would conclude
  that SAXPY enjoys the full $\approx6\times$ SIMD speedup. We predict instead that the same
  \emph{vectorized} kernel shows almost no speedup on a DRAM-resident array, because memory bandwidth,
  not vector throughput, then sets the pace.
\end{hypothesis}

\begin{knobs}
  Case        & memory-bound SAXPY: a loop that provably vectorizes (\code{vfmadd213ps} on
                \code{ymm}, \cref{sec:vectorization}) but whose working set far exceeds L3 \\
  Build       & \code{native} (\code{-O3 -march=native}) vs \code{scalar} \\
  Script      & \code{SIMD_Advantage_Profiling/scripts/run.sh negative} \\
\end{knobs}

\subsection{Evidence}
The SAXPY loop at \code{kernels.c:29} is vectorized at \qty{256}{bit} with FMA
(\cref{sec:vectorization}), yet its speedup over scalar falls with size:
\begin{center}\small
  \begin{tabular}{@{}rrrrr@{}}
    \toprule
    $N$ & working set & native GFLOP/s & \% of roof & speedup \\
    \midrule
    \num{1024}      & \qty{8}{\kibi\byte} (L1)    & 20.1 & 72  & $5.60\times$ \\
    \num{16384}     & \qty{128}{\kibi\byte} (L2)  & 13.7 & 84  & $4.11\times$ \\
    \num{262144}    & \qty{2}{\mebi\byte} (L3)    & 8.6  & 99  & $2.82\times$ \\
    \num{4194304}   & \qty{32}{\mebi\byte} (DRAM) & 6.0  & --  & $2.31\times$ \\
    \num{8388608}   & \qty{64}{\mebi\byte} (DRAM) & 5.2  & 120 & $2.19\times$ \\
    \num{16777216}  & \qty{128}{\mebi\byte} (DRAM) & 4.7 & 109 & $1.95\times$ \\
    \num{33554432}  & \qty{256}{\mebi\byte} (DRAM) & 4.5 & 104 & $1.92\times$ \\
    \bottomrule
  \end{tabular}
\end{center}
``\% of roof'': native GFLOP/s over $B_\ell \cdot I$ for the level holding the data
(\cref{sec:roofline}); \qty{32}{\mebi\byte} straddles L3 and DRAM.

\subsection{Explanation}
The hypothesis was half right: the $\approx6\times$ of L1 does not survive in DRAM, but the speedup
settles at $1.92\text{--}1.95\times$, not ``almost none''. The loop genuinely uses SIMD, but at
\qty{0.167}{\flop\per\byte} it is deep in the bandwidth-bound region: native SAXPY runs at \SIrange{104}{120}{\percent} of the separately measured
DRAM roof with the compulsory \numrange{0.122}{0.128} LLC misses per element (\cref{sec:locality}).
\emph{No} wider vector, AVX-512 included, could make it faster. The remaining $\approx1.9\times$ is
the scalar loop's instruction cost (\num{2.34} vs \num{1.25} cycles per element), which shrinks at
turbo clocks (\cref{sec:method}).

The dot product is subtler: it reports ``vectorized'' and uses a \qty{256}{bit} \code{vmulps},
yet the serial reduction leaves $\approx1\times$. Both carry the same lesson for AI-assisted benchmarking: ``it compiled to SIMD'' is necessary but not
sufficient evidence that a kernel is faster.

% !TeX root = ../main.tex
\section{Anomalies and Limitations}\label{sec:anomalies}
\rubric{Report quality}{10}{Clear plots, concise writing, honest anomalies and thoughtful limitations.}

\subsection{Unexpected results}
The dot product, the DRAM gap and the losing AVX2 gather are explained in
\cref{sec:vectorization,sec:locality,sec:stride}. Three anomalies remain open:
\begin{itemize}
  \item \textbf{Cache-resident stencil is slower than L3-resident stencil}:
        \qtyrange{16.6}{17.3}{\giga\flop\per\second} for $N=\num{1024}\ldots\num{16384}$ but
        \qtyrange{19.9}{20.7}{\giga\flop\per\second} for $N=\num{65536}\ldots\num{1048576}$,
        reproducibly, which no monotonic cache effect explains. Untested hypothesis: 4K aliasing
        between stores to \code{y[i]} and loads of \code{x[i+1]}; sweeping the \code{x}--\code{y}
        distance modulo \qty{4}{\kibi\byte} with the \code{ld_blocks_partial.address_alias} counter
        would test it.
  \item \textbf{Native SAXPY peaks within L1} ($6.86\times$ at $N=\num{1536}$, $6.32\times$ at
        $N=512$, else $5.6\text{--}5.9\times$); the same placement hypothesis applies.
  \item \textbf{Same code, different process, $N=\num{32768}$}: \SI{1}{\percent} (scalar) and
        \SI{14}{\percent} (native) apart. Untested hypothesis: the two arrays exactly fill the
        \qty{256}{\kibi\byte} L2, so conflict misses depend on physical page placement.
\end{itemize}

\subsection{Limitations}
\begin{itemize}
  \item \textbf{Core-PMU counters only}: no DRAM-traffic or TLB events, and the generic LLC-miss
        event's prefetch handling is model-specific, so we only compare it between builds. IMC bytes
        per element (as read by the companion Cache \& Memory report) would turn the roofline's
        traffic assumption into a measurement.
  \item \textbf{Fixed clock}: reproducible but not normal operation; at turbo, DRAM-resident SAXPY
        speedups shrink to $1.32\text{--}1.44\times$ (\cref{sec:method}).
  \item \textbf{One machine, one compiler} (i9-10900F, GCC \sysGCC{}), no AVX-512 part.
  \item \textbf{THP always-on, no core isolation, unshuffled sweep order}; the shared-core checks of
        \cref{sec:method} guard the data.
  \item \textbf{Discarded sessions}: earlier sessions that day, before the GPE storm was masked and
        while another agent's unlocked smoke tests briefly shared CPU~2.
\end{itemize}

% !TeX root = ../main.tex
\section{AI-Use Note}\label{sec:ai-use}

Summarized from \code{AI_USE_LOG.md} in the code repository.

\begin{description}[style=nextline]
  \item[Tools used] Claude Code, an AI coding agent, for coding, debugging, and running the experiments.
        All sessions used the model Claude Opus~5.5. On
        2026-09-15 it produced the templates (folder structure, READMEs, Makefiles,
        files with placeholder markers, report skeletons) that every later step filled in. On 2026-10-12 to
        2026-10-20 it produced the implementation, and on 2026-10-21 the rework (measured ceilings,
        counters, AVX2 gather, extended sweeps and the fixed-clock re-run). It also produced a
        preliminary draft of the report, which I discarded; I wrote this report myself.
        GCC~\sysGCC{} and \code{objdump} produced the builds and the vectorization evidence.
  \item[Who ran the measurements] The timed SIMD experiments of 2026-10-21 (the reported runs, plus earlier sessions that
        were discarded for interference) were executed by Claude
        Code on my machine, with my approval. Before the runs I set the
        governor to \code{performance}, disabled turbo and SMT, and masked the ACPI GPE~0x42
        interrupt storm. Claude Code checked these settings before each run and serialized its runs with
        other benchmark jobs through a lock file. All measurements come from these real runs on
        the machine in \cref{tab:environment}.
  \item[What they helped generate] The four kernels and their references; the build-variant
        Makefile; the benchmark driver and the shared timer/allocation/pinning helpers; the
        experiment scripts (\code{run.sh}), the summarizer (\code{summarize.c}) and the
        vectorization-evidence script; the roofline-ceiling microbenchmarks
        (\code{ceilings.c}), the \code{perf_event_open} counter helper (\code{perfctr.c}), the
        explicit AVX2 gather kernel and the post-processing scripts (\code{speedup_iqr.awk},
        \code{roofline.sh}, \code{tables.sh}).
  \item[Who verified what] The AI ran the automated checks: \code{make test} (\num{578} cases in total),
        the disassembly comparisons, the per-run settings and shared-core checks, and the
        CSV-to-text spot checks of every quoted number. I verified all AI output, applied
        the root-only system settings, and reviewed the results,
        which drove several of the corrections below. ``Independently'' therefore means checked
        against measurements and tool output, as well as against my own analysis.
  \item[What was independently modified, tested, or rejected] Concrete corrections, among others.
        (1)~The generated Makefile compiled the \emph{whole} program with each variant's
        flags; this was restructured so only the kernel translation unit varies, keeping the
        timing harness identical across builds. (2)~The processed-CSV column names and one
        figure filter had drifted from what the plots expect---the stride figure asked for
        \code{n=16777216} while the data was measured at \code{n=4194304}---which would have
        produced a silently empty plot; the filter was corrected to the measured size.
        (3)~\emph{An AI-written script was simply wrong}: the \code{awk} in
        \code{verify_vectorization.sh} took its output path as a trailing \code{VAR=value}
        argument, which \code{awk} does not assign until it reaches that argument, so the
        \code{BEGIN} block saw an empty string and aborted with a null-redirection error. It
        was caught the first time \code{make report} was run and fixed by passing the path
        with \code{-v}. We also chose \emph{not} to enable \code{-ffast-math}, rejecting the
        easy way to make the dot product vectorize, because it would break bit-reproducibility
        against the reference---and because the non-reassociated reduction is itself a result
        worth reporting (\cref{sec:vectorization,sec:negative}).
  \item[Validation] Correctness: \num{578} reference and known-answer cases in total
        (\num{190} scalar, \num{190} autovec, \num{198} native) pass (\cref{sec:correctness}).
        Vectorization: confirmed by compiler report and disassembly, not assumed
        (\cref{sec:vectorization}). Numbers: every number from the 2026-10-21 runs traces to a
        committed file (\cref{app:repro}); the memory-bound claims are
        cross-checked against bandwidth and peak FLOP rate measured by separate microbenchmarks
        (\code{src/ceilings.c}), not derived from the kernels themselves (\cref{sec:roofline}).
        Review corrections that changed results: an early roofline bandwidth was derived from
        SAXPY itself (circular; replaced by the independent ceilings); the AI-drafted claims that
        scalar and SIMD ``converge'' in DRAM and that small-$N$ speedups are inflated by per-call
        overhead were both contradicted by the fixed-clock data and ignored; the alignment sweep had
        silently skipped offsets \qtyrange{40}{56}{\byte} and was re-run; and the L1 roof first
        included a write-allocate read that L1-resident stores do not incur.
\end{description}

% !TeX root = ../main.tex
\section{Conclusion}\label{sec:conclusion}

On this AVX2 machine SIMD pays off handsomely only when the data fits in cache and the kernel is
not a serial reduction: \code{-march=native} then gives $3.5$--$5.8\times$ at $N=\num{4096}$,
limited by load/store issue rather than the FMA units. In DRAM the vector builds run \emph{at} the independently
measured bandwidth roof, and what remains ($1.88$--$2.02\times$ for SAXPY) is the scalar loops' own
instruction cost at the fixed clock. Float64 keeps about half the cache-resident gain, and any
non-unit stride or gather erases it. The main lesson: ``the loop vectorized'' is not ``the loop
is fast''.

% % !TeX root = ../main.tex
% Optional CPU-vs-GPU extension. Enable with % !TeX root = ../main.tex
% Optional CPU-vs-GPU extension. Enable with % !TeX root = ../main.tex
% Optional CPU-vs-GPU extension. Enable with \input{sections/15_gpu_extension} in main.tex.
\section{Optional Extension: CPU vs.\ GPU}\label{sec:gpu}

GPU: \sysGPU; CUDA \sysCUDA.

\begin{hypothesis}
  \todo{Break-even $N$ you expect, once host--device transfer is included.}
\end{hypothesis}

\begin{knobs}
  Kernel      & \tbd \\
  Timing      & kernel-only (device synchronized) and end-to-end (including host$\leftrightarrow$device copies) \\
  Varied      & $N$ \\
\end{knobs}

\todo{Figure: CPU (best SIMD build) vs.\ GPU kernel-only vs.\ GPU end-to-end throughput across $N$.}

\subsection{Analysis}
\todo{When GPU acceleration helps, when transfers erase the benefit, and why.}
 in main.tex.
\section{Optional Extension: CPU vs.\ GPU}\label{sec:gpu}

GPU: \sysGPU; CUDA \sysCUDA.

\begin{hypothesis}
  \todo{Break-even $N$ you expect, once host--device transfer is included.}
\end{hypothesis}

\begin{knobs}
  Kernel      & \tbd \\
  Timing      & kernel-only (device synchronized) and end-to-end (including host$\leftrightarrow$device copies) \\
  Varied      & $N$ \\
\end{knobs}

\todo{Figure: CPU (best SIMD build) vs.\ GPU kernel-only vs.\ GPU end-to-end throughput across $N$.}

\subsection{Analysis}
\todo{When GPU acceleration helps, when transfers erase the benefit, and why.}
 in main.tex.
\section{Optional Extension: CPU vs.\ GPU}\label{sec:gpu}

GPU: \sysGPU; CUDA \sysCUDA.

\begin{hypothesis}
  \todo{Break-even $N$ you expect, once host--device transfer is included.}
\end{hypothesis}

\begin{knobs}
  Kernel      & \tbd \\
  Timing      & kernel-only (device synchronized) and end-to-end (including host$\leftrightarrow$device copies) \\
  Varied      & $N$ \\
\end{knobs}

\todo{Figure: CPU (best SIMD build) vs.\ GPU kernel-only vs.\ GPU end-to-end throughput across $N$.}

\subsection{Analysis}
\todo{When GPU acceleration helps, when transfers erase the benefit, and why.}
   % optional CPU-vs-GPU comparison

\printbibliography

\appendix
% !TeX root = ../main.tex
\section{Reproducibility}\label{app:repro}

\newcommand{\acscommit}{\detokenize{d1e7b4b7fc69}}
Code and data (\code{code/}) and this report (\code{latex/}) live in the project's folder of one repository, with shared files in \code{common/}. The timed runs (2026-10-21, 20:38--20:54) used the tree committed as \texttt{b00db66}; the
measurement code is unchanged since, except an opt-in error print in the validation path of
\code{main.c} that is inactive in timed runs. This report corresponds to commit
\texttt{\acscommit}. From the repository root:

\begin{lstlisting}[language=bash, basicstyle=\ttfamily\scriptsize]
gh repo clone SKempiak/ACS-Code ACS && cd ACS
git checkout d1e7b4b7fc69
make tools                                     # common/code/build/tools/summarize
make -C SIMD_Advantage_Profiling/code          # scalar/autovec/native simd_bench + ceilings
make -C SIMD_Advantage_Profiling/code test     # 578 correctness cases (max error logged)
make -C SIMD_Advantage_Profiling/code report   # vectorization evidence
common/code/scripts/collect_env.sh --tex > common/latex/system_info.tex
SIMD_Advantage_Profiling/code/scripts/run.sh all     # ceilings, baseline, locality, alignment,
                                                     # stride, dtype, negative
common/code/scripts/sync_data.sh SIMD_Advantage_Profiling   # -> latex/data/
make simd                                      # -> latex/SIMD_Advantage_Profiling.pdf
\end{lstlisting}
Before timing, as root: governor \code{performance}, \code{intel_pstate/no_turbo}~$=1$, SMT off,
and \code{echo disable > /sys/firmware/acpi/interrupts/gpe42} (none persists across a reboot).
\code{run.sh derive} recomputes the derived tables without measuring. Each raw directory
(\cref{tab:raw-data}) holds \code{raw.csv} and \code{env.txt}; the turbo comparison traces to
commit \code{73f8fb8}.

\begin{table}[h]
  \centering\footnotesize
  \caption{Raw data behind each figure and table.}\label{tab:raw-data}
  \begin{tabularx}{\linewidth}{@{}lX@{}}
    \toprule
    Figure / table & Raw run directory (under \code{SIMD_Advantage_Profiling/code/data/raw/}) \\
    \midrule
    \Cref{fig:baseline-speedup,fig:baseline-throughput} & \code{baseline/20261021-203854} \\
    \Cref{fig:locality,fig:locality-bw,fig:locality-counters}, \cref{tab:locality-bw} & \code{locality/20261021-204801} \\
    \Cref{fig:alignment,fig:tail} & \code{alignment/20261021-204858} (\code{raw.csv}, \code{raw_tail.csv}) \\
    \Cref{fig:stride}, \cref{tab:stride} & \code{stride/20261021-205143} \\
    \Cref{fig:dtype}, \cref{tab:dtype} & \code{dtype/20261021-205209} \\
    \Cref{tab:baseline-gflops} & \code{baseline/20261021-203854} \\
    Correctness result (\cref{sec:correctness}) & \code{validation/20261021-231629} (\code{make_test.txt}) \\
    \Cref{fig:roofline}, \cref{tab:ceilings} & \code{ceilings/20261021-204708}, \code{baseline/20261021-203854} \\
    Negative-case table (\cref{sec:negative}) & \code{negative/20261021-205359} \\
    \Cref{tab:vec-summary} & \code{vectorization/20261021-001032} \\
    \bottomrule
  \end{tabularx}
\end{table}


\end{document}
